% The vacuum Bell correlation of the massive free scalar field attains the
% classical bound for bounded regions beyond a finite separation.
% Physical Superintelligence (PSI), draft of 22 October 2026.
% Compile with: pdflatex oqp12_prl_draft; bibtex oqp12_prl_draft; pdflatex; pdflatex

\documentclass[aps,prl,reprint,longbibliography,raggedbottom]{revtex4-2}
\usepackage[T1]{fontenc}
\usepackage{amsmath,amssymb,amsthm}
\usepackage{newtxtext,newtxmath}
\usepackage{microtype}
\usepackage{booktabs}
\PassOptionsToPackage{hyphens}{url}
\usepackage[hidelinks]{hyperref}
\makeatletter\g@addto@macro{\UrlBreaks}{\do\-}\makeatother
\hypersetup{
  pdftitle={The vacuum Bell correlation of the massive free scalar field attains the classical bound for bounded regions beyond a finite separation},
  pdfauthor={David Turturean}
}
\allowdisplaybreaks[1]
\setcounter{secnumdepth}{3}

\newtheorem{theorem}{Theorem}
\newtheorem{lemma}{Lemma}
\newtheorem{corollary}{Corollary}

\newcommand{\C}{\mathbb C}
\newcommand{\R}{\mathbb R}
\newcommand{\N}{\mathbb N}
\newcommand{\ii}{\mathrm i}
\newcommand{\dd}{\,\mathrm d}
\newcommand{\norm}[1]{\lVert #1\rVert}
\newcommand{\abs}[1]{\lvert #1\rvert}
\newcommand{\ip}[2]{\langle #1,#2\rangle}
\DeclareMathOperator{\Tr}{Tr}
\DeclareMathOperator{\Rea}{Re}
\DeclareMathOperator{\Ima}{Im}
\DeclareMathOperator{\supp}{supp}
\DeclareMathOperator{\rank}{rank}
\DeclareMathOperator{\diag}{diag}
\DeclareMathOperator{\vol}{vol}
\DeclareMathOperator{\Ran}{Ran}
\DeclareMathOperator{\Dom}{Dom}
\DeclareMathOperator{\artanh}{artanh}
\newcommand{\A}{\mathfrak A}
\newcommand{\B}{\mathfrak B}
\newcommand{\Hh}{\mathcal H}
\newcommand{\Res}{\mathsf R}
\newcommand{\JointRes}{\Res_{U_s}}
\newcommand{\JointD}{\mathsf D}
\newcommand{\Therm}{\mathsf H}
\newcommand{\Qtil}{\widetilde Q}
\newcommand{\cov}{\mathsf v}
\newcommand{\sx}{\mathsf f}
\newcommand{\sy}{\mathsf f'}
\newcommand{\bx}{\mathbf x}
\newcommand{\by}{\mathbf y}
\newcommand{\bp}{\mathbf p}
\newcommand{\bz}{\mathbf z}
\newcommand{\eone}{\mathbf e_1}
\newcommand{\Energy}{\mathcal E}
\newcommand{\Weyl}{\mathsf W}
\providecommand{\doi}[1]{\href{https://doi.org/#1}{\nolinkurl{doi:#1}}}
\providecommand{\arxiv}[1]{\href{https://arxiv.org/abs/#1}{\nolinkurl{arXiv:#1}}}

\AtBeginDocument{\providecommand{\doi}{}\renewcommand{\doi}[1]{\href{https://doi.org/#1}{\nolinkurl{doi:#1}}}}
\begin{document}

\title{The vacuum Bell correlation of the massive free scalar field
attains the classical bound for bounded regions beyond a finite
separation}
\author{David Turturean}\thanks{D.T. built and ran the automated research and formalization workflows at PSI that produced these results, steered the writing of the paper, and asked for AI reviews and revisions of the drafts.}
\affiliation{Physical Superintelligence PBC, Cambridge, Massachusetts, USA}

\begin{abstract}
For the real free scalar field in $3+1$ dimensions we prove that a
universal finite $C>2$ exists: for every mass $m>0$ and radius $R>0$,
the vacuum Bell correlation of two equal-time double cones of radius
$R$ equals its classical bound at centre separations $L\ge CR$. The
supremum defining the Bell correlation includes all self-adjoint local
contractions; the vacuum remains entangled. By isotony, CHSH violation
is absent for fixed bounded regions moved sufficiently far apart along
a fixed spacelike direction. This resolves the fixed-bounded-region
version of Open Quantum Problem~12.
\end{abstract}

\maketitle

{\raggedright\setlength{\parskip}{0.4\baselineskip}
\section*{Acknowledgements}

We thank Adam Levine for condensing the full manuscript into this letter and for his comments on the presentation, and Ning Bao for a referee report on an earlier draft of this letter; Ning Bao wrote the report as a paid contractor to PSI, not under one of his grants or as an employee of Northeastern University. We also thank Physical Superintelligence PBC, in which the author has a financial interest, for use of its resources, including its technology and frontier AI models. The author used these models to develop, derive, and compose this manuscript and its Lean~4 supplements.

Results were formalized in Lean~4 by AI agents and checked by Lean's kernel; the formalization is at \url{https://github.com/physical-superintelligence/oqp-batch-2026/tree/main/oqp12-vacuum-bell-correlations/lean}.

David Turturean reviewed and edited this paper and is responsible for its correctness. David Turturean is employed by PSI and has a financial interest in it.

\par}
\par\medskip\noindent\rule{\columnwidth}{0.4pt}\par\medskip


\section{Statement of the result}\label{sec:statement}

Let $O_1,O_2$ be spacelike separated regions of Minkowski space, with
local von Neumann algebras $\A=\A(O_1)$ and $\B=\A(O_2)$, which commute,
and let $\omega$ be a state. An observable measurable in $O_1$ whose
values lie in $[-1,1]$ is a self-adjoint $A\in\A$ with $\norm A\le1$,
a self-adjoint \emph{contraction} ($\norm\cdot$ is the operator norm).
The CHSH inequality \cite{CHSH69} concerns two such observables on each
side. Following Summers and Werner \cite{SW87a}, the \emph{Bell
correlation} of $\omega$ is
\begin{equation}
\beta(\omega;\A,\B)
=\sup\tfrac12\bigl\lvert\omega\bigl(A_0(B_0+B_1)+A_1(B_0-B_1)\bigr)\bigr\rvert,
\label{eq:beta}
\end{equation}
the supremum over self-adjoint $A_i\in\A$ and $B_j\in\B$ of norm at
most $1$. The CHSH inequality is the statement that in a \emph{local
hidden variable model}, in which the four measured values are
functions $a_i(\lambda),b_j(\lambda)\in[-1,1]$ of a shared classical
variable $\lambda$ and expectations are averages over $\lambda$,
\begin{equation}
\tfrac12\abs{\langle a_0b_0\rangle+\langle a_0b_1\rangle
+\langle a_1b_0\rangle-\langle a_1b_1\rangle}\le1 ,
\label{eq:CHSH}
\end{equation}
because $\abs{a_0(b_0+b_1)+a_1(b_0-b_1)}\le\abs{b_0+b_1}+\abs{b_0-b_1}\le2$
for every $\lambda$. Quantum mechanics allows values of $\beta$ up to $\sqrt2$~\cite{Tsirelson80}:
for \emph{reflections}, observables with $A_i^2=I$
and $B_j^2=I$ and hence with outcomes $\pm1$, one has
${\bigl(A_0(B_0+B_1)+A_1(B_0-B_1)\bigr)^2}=4I-[A_0,A_1][B_0,B_1]$, of norm
at most $8$. The choice $A_i=B_j=I$ gives $1$, so $\beta\ge1$ always.
If $\beta>1$, the correlations of $\omega$ between the two regions
cannot be reproduced by any local hidden variable model; this is a
\emph{Bell violation}. Note that the supremum in \eqref{eq:beta} is unchanged if
it is restricted to reflections, because every self-adjoint contraction is an
average of reflections and the expression is affine in each argument (Appendix~\ref{app:probability}, Step~1). (Verch and Werner
\cite{VW05} omit the factor $\frac12$, so that their bounds are $2$ and
$2\sqrt2$.)

By the Reeh--Schlieder theorem the vacuum is entangled across any two
spacelike separated regions \cite{HC00}, and the entanglement is
distillable at any distance \cite{VW05}: local operations and
classical communication can extract standard entangled pairs from it.
Bell correlation is a
stronger property. Summers and Werner proved $\beta=\sqrt2$ for the
vacuum of a massive free field on complementary wedges
($\{x^1>\abs{x^0}\}$ and its spatial mirror) and on tangent
double cones \cite{SW85,SW87b,SW87c,SW88}, and for regions at a
positive distance a mass gap $m>0$ gives
$\beta\le1+2e^{-m\,t_{12}}$, where $t_{12}$ is the largest time
through which one region can be shifted while the two stay spacelike
separated \cite{SW87a,Summers11}. This is an upper bound; whether some
violation survives at every distance is Problem~12 of the IQOQI Vienna
Open Quantum Problems list~\cite{OQP12,KW05}: ``Decide whether some
(necessarily small)
violation of the inequalities is possible for regions arbitrarily far
apart. For definiteness, consider a massive scalar free relativistic
Bose field.'' Violations obtained after local filtering of coupled
probes \cite{RRS05,MN20} concern a post-selected state and are not at
issue here. Nor is entanglement itself: suitable local modes
detect vacuum entanglement at every distance \cite{GK25}, but
entanglement of the unfiltered vacuum does not imply a CHSH
violation.

\emph{Setting.---}Units $\hbar=c=1$, signature $(-,+,+,+)$, so that
$x\cdot y=-x^0y^0+\bx\cdot\by$ for points $(x^0,\bx)$ with $\bx\in\R^3$;
$B_R(\bx_0)=\{\bx\in\R^3:\abs{\bx-\bx_0}<R\}$ is the open ball of radius
$R$; Fourier transforms are
$\hat f(\bp)=\int e^{-\ii\bp\cdot\bx}f(\bx)\dd^3x$, and
$\eone=(1,0,0)$. For $R>0$ the double cone
$D_R=\{(x^0,\bx):\abs{x^0}+\abs\bx<R\}$ has base $B_R(0)$ at time zero,
and $D_R+(0,L\eone)$ has its centre at distance $L$ along the first
axis. We consider the real free scalar field of mass $m>0$ in $3+1$
dimensions, its vacuum state $\omega$, and its local algebras
$\A_m(O)$. For $L>2R$ the two bases are disjoint, the algebras of the
two double cones commute, and we write
\[
\beta_m(R,L)=\beta\bigl(\omega;\A_m(D_R),\A_m(D_R+(0,L\eone))\bigr).
\]
The supremum in \eqref{eq:beta} runs over all bounded observables of
the two regions; they may depend on $L$ and may involve events of
arbitrarily small probability.

\begin{theorem}[Finite separation cutoff]\label{thm:main}
There is a finite constant $C_*>2$, independent of $m$ and $R$, such
that
\begin{equation}
\forall\,m,R>0,\qquad
L\ge C_*R\quad\Longrightarrow\quad\beta_m(R,L)=1.
\label{eq:main}
\end{equation}
\end{theorem}

\begin{corollary}[Fixed bounded regions]\label{cor:fixed}
For every $m>0$, any two bounded regions $O_1,O_2$ and any fixed
spacelike direction, the vacuum Bell correlation of $\A_m(O_1)$ and of
the local algebra of the translate of $O_2$ equals $1$ for all
sufficiently large translations of $O_2$ in that direction.
\end{corollary}

This answers Problem~12 in the negative for fixed bounded
regions. The theorem supplies a finite sufficient separation, not a
sharp threshold: it does not locate the smallest such separation, and
for regions that grow with the separation, and for wedges at a
positive gap, the question remains open. The theorem is an existence
statement, and the constant $C_*$
it produces is finite but not explicit. What the theorem excludes is a CHSH
violation, however small and however rare its outcomes, between fixed
bounded regions far enough apart.

The proof compares the vacuum with the product of its
restrictions to the two regions: the state with the same local physics
and no correlations between the regions at all.
Section~\ref{sec:probability} turns closeness to this product state
into the classical bound $\beta=1$, and Sections \ref{sec:gaussian}
and~\ref{sec:estimates} establish the closeness at large separation,
uniformly in the mass.

Section~\ref{sec:field} sets up the field and makes two reductions;
Sections~\ref{sec:probability}--\ref{sec:completion} prove
Theorem~\ref{thm:main}.


\section{Setting and two reductions}\label{sec:field}

At time zero the field is described by the smeared field and momentum
\[
\phi(f)=\int\phi(0,\bx)f(\bx)\dd^3x,\qquad
\pi(g)=\int\dot\phi(0,\bx)g(\bx)\dd^3x,
\]
$f,g\in C_c^\infty(\R^3;\R)$, where
$[\phi(0,\bx),\dot\phi(0,\by)]=\ii\delta^3(\bx-\by)$. We call $\sx=(f,g)$ a
\emph{source}, denoted by the sans-serif letter of its field component,
and write 
\begin{align}\label{eq:comm}
\Phi(\sx)=\phi(f)+\pi(g)
\end{align} and
$\Weyl(\sx)=e^{\ii\Phi(\sx)}$. The commutator is
$[\Phi(\sx),\Phi(\sy)]=\ii\varsigma(\sx,\sy)$ with
\begin{equation}
\varsigma(\sx,\sy)=\int_{\R^3}(fg'-gf')\dd^3x,\qquad \sy=(f',g'),
\label{eq:symplectic}
\end{equation}
so $\Weyl(\sx)\Weyl(\sy)=e^{-\ii\varsigma(\sx,\sy)/2}\Weyl(\sx+\sy)$, and
$\Weyl(\sx)$ and $\Weyl(\sy)$ commute when the supports of $\sx$ and
$\sy$ are disjoint ($\supp\sx:=\supp f\cup\supp g$), since then
$\varsigma(\sx,\sy)=\varsigma(\sy,\sx)=0$. For mass
$\mu>0$, the vacuum is the
Gaussian (quasi-free) state \cite{Haag96}
\begin{equation}
\begin{gathered}
\omega(\Weyl(\sx))=e^{-\cov_\mu(\sx,\sx)/2},\\
\cov_\mu(\sx,\sy)=\tfrac12\bigl(\ip f{(\mu^2-\nabla^2)^{-1/2}f'}+\ip g{(\mu^2-\nabla^2)^{1/2}g'}\bigr),
\end{gathered}
\label{eq:covariance}
\end{equation}
with $\ip\cdot\cdot$ the $L^2(\R^3)$ inner product; $\cov_\mu$ is the
\emph{covariance}, and $\cov_\mu(\sx,\sx)=\omega(\Phi(\sx)^2)$. (The
covariances of \cite{SLW18} are twice ours.) By Araki's time-slice
theorem \cite{Araki64,EO73},
\begin{equation}
\A_\mu(D_R)=\{\Weyl(f,g):f,g\in C_c^\infty(B_R(0);\R)\}'' ,
\label{eq:time-slice}
\end{equation}
(the double prime is the double commutant, the smallest von Neumann
algebra containing the set in braces), and likewise for the translated
double cone with sources in
$B_R(L\eone)$. An \emph{effect} is an operator $0\le E\le I$, the
quantum analogue of an event, with probability $\omega(E)$; a
\emph{projection} is an effect with $E^2=E$; a self-adjoint
contraction $A$ has the two effects $(I\pm A)/2$, which are projections
when $A^2=I$, in which case we call $A$ a \emph{reflection}.

\begin{lemma}[Dilation; cf.~{\cite[Section~2]{BV98}}]\label{lem:scaling}
$\beta_m(R,L)=\beta_{mR}(1,L/R)$ for all $m,R>0$ and $L>2R$.
\end{lemma}

\begin{proof}
The map $(f,g)\mapsto(\tilde f,\tilde g)$, $\tilde f(\by)=R^2f(R\by)$,
$\tilde g(\by)=Rg(R\by)$, carries sources in $B_R(0)$ and $B_R(L\eone)$
to sources in $B_1(0)$ and $B_1(s\eone)$, $s=L/R$, preserves
\eqref{eq:symplectic}, and carries the covariance of mass $m$ to that
of mass $mR$, since $\hat{\tilde f}(\mathbf q)=R^{-1}\hat f(\mathbf q/R)$,
$\hat{\tilde g}(\mathbf q)=R^{-2}\hat g(\mathbf q/R)$ and
$\sqrt{(mR)^2+\abs{\mathbf q}^2}=R\sqrt{m^2+\abs{\mathbf q/R}^2}$. It therefore induces a unitary
between the two vacuum representations that fixes the vacuum and maps
each local algebra of \eqref{eq:time-slice} onto its counterpart, so the
two suprema \eqref{eq:beta} coincide.
\end{proof}

From now on $\mu>0$ is arbitrary (the letter $\mu$ keeps this mass
distinct from the physical $m$ of Theorem~\ref{thm:main}), the balls are $B_1(0)$ and
$B_1(s\eone)$ with $s>2$, and
\[
\A=\A_\mu(D_1),\qquad \B=\A_\mu(D_1+(0,s\eone)).
\]
Theorem~\ref{thm:main} is the statement that $\beta(\omega;\A,\B)=1$
whenever $s\ge C_*$.

\emph{One-particle vectors.---}The one-particle space of the field of
mass $\mu$ is $L^2(H_\mu^+,\dd\mathfrak m_\mu)$,
$H_\mu^+=\{p:p_0=\sqrt{\mu^2+\abs\bp^2}\}$,
$\dd\mathfrak m_\mu=\dd^3p/(2(2\pi)^3p_0)$, and the one-particle vector
of a source is
\begin{equation}
(\kappa\sx)(p)=\hat f(\bp)+\ii p_0\hat g(\bp) .
\label{eq:kappa}
\end{equation}
By Plancherel's formula, for real sources,
\begin{equation}
\ip{\kappa\sx}{\kappa\sy}=\cov_\mu(\sx,\sy)+\tfrac\ii2\varsigma(\sx,\sy)
\label{eq:gram}
\end{equation}
(inner products are antilinear in the first argument; on the left it
is that of $L^2(H_\mu^+,\dd\mathfrak m_\mu)$).
For a nonempty bounded open $U\subset\R^3$, let $\mathcal X_{\mu,U}$ be
the completion of the real smooth sources in $U$ in the norm
$\cov_\mu(\sx,\sx)^{1/2}$, and $\Hh_{\mu,U}$ its complexification, with inner product
$\ip\cdot\cdot_{\mu,U}$ extending $\cov_\mu$ sesquilinearly; by
\eqref{eq:gram}, $\kappa$ extends to a bounded map
$\kappa_U\colon\Hh_{\mu,U}\to L^2(H_\mu^+,\dd\mathfrak m_\mu)$,
isometric on real sources. Real-linear maps such as inclusions of
source spaces are extended complex-linearly without change of
notation.

\begin{lemma}[Injectivity and density]\label{lem:standard}
$\kappa_U$ is injective and has dense range.
\end{lemma}

The proof (Appendix~\ref{app:evaluation}) uses that the Fourier
transforms of compactly supported sources are entire. The statement that $\kappa_U$ has dense range is the one-particle form of the Reeh--Schlieder
property; injectivity is what will give the strict uncertainty relation
\eqref{eq:strict-uncertainty} below.

\subsection{Modes and finite canonical systems}
A \emph{mode} is one
canonical degree of freedom of the field in a region. This corresponds to a pair of sources
$\sx,\sy$ there with $\varsigma(\sx,\sy)=1$, so that
$[\Phi(\sx),\Phi(\sy)]=\ii$ by \eqref{eq:comm}. The local algebras of a free massive scalar carry no density matrices
(Section~\ref{sec:gaussian}), so every comparison in this paper is
made on finitely many modes at a time: the degrees of freedom seen by
a detector that couples to finitely many smeared fields $\Phi(\sx)$ in
the region.

A \emph{finite canonical system} is a finite collection of
independent modes. Formally, it is a
finite-dimensional real space $F$
of smooth sources in a bounded open set, on which $\varsigma$ is
nondegenerate: in a basis with
$\varsigma(\sx_k,\sx_l)=(J_n)_{kl}$, $J_n=\bigoplus_{j=1}^n\left(\!\begin{smallmatrix}0&1\\-1&0\end{smallmatrix}\!\right)$,
the $2n$ operators $\Phi(\sx_k)$ form $n$ canonical pairs,
$[\Phi(\sx_{2j-1}),\Phi(\sx_{2j})]=\ii$ with all other commutators
zero. 

By the Stone--von Neumann theorem \cite{Haag96}, the algebra generated by
$\Weyl(\sx)$, $\sx\in F$, is a copy of the bounded operators on
$L^2(\R^n)$---ordinary quantum mechanics of $n$ degrees of
freedom---possibly with multiplicity, and the vacuum restricted to
it is given by a density matrix, since a vector state is normal.
By \eqref{eq:covariance} its characteristic function is Gaussian,
\begin{align}\label{eq:Weyl}
\omega\Bigl(\Weyl\Bigl(\sum_k\xi_k\sx_k\Bigr)\Bigr)=e^{-\xi^TV\xi/2}
\end{align}
for $\xi\in\R^{2n}$, and a density matrix is determined by its
characteristic function \cite{SLW18}; so the restriction is the
Gaussian state $\rho_F$ with
covariance matrix $V_{kl}=\cov_\mu(\sx_k,\sx_l)$.

A finite collection of sources in an open set $U$ need not form
canonical pairs by itself: field sources $(f,0)$, for instance, all
commute. Nevertheless, such a finite collection can always be completed into a finite canonical system
supported inside $U$ by adjoining finitely many further sources. The only
obstruction to such a construction would be if there were a direction $\sx=(f,g)$ in the span that commutes with
the whole span, but we can always include the missing partner $(0,f)$ since $\varsigma(\sx,(0,f))=\int f^2>0$
if $f\ne0$ or $(-g,0)$ if $f=0$ since $\varsigma(\sx,(-g,0))=\int g^2>0$. Each such adjoined
partner lowers the number of independent partnerless
sources by one, so finitely many steps complete the pairs. Such a construction allows us to always complete
two finite canonical systems in $U$ into a common third system,
which Lemma~\ref{lem:exhaustion} below needs. 

Given a Gaussian state on a given finite canonical system, Williamson's normal
form \cite{Williamson36} writes $V=S(\bigoplus_j\upsilon_jI_2)S^T$ with
$S$ a linear canonical transformation and $I_2$ the identity on the corresponding canonical pair. In these normal-mode
coordinates $\rho_F$ is a product of $n$ oscillator states, the $j$th
with covariance $\upsilon_jI_2$. Each normal mode is thermal:
with $N_j$ its number operator, the $j$th factor is proportional to
$e^{-\beta_jN_j}$ for an inverse temperature $\beta_j\in(0,\infty]$,
and the \emph{symplectic eigenvalue} $\upsilon_j\ge\frac12$ is related to the inverse temperature of the mode by 
\begin{equation}
\upsilon_j=\tfrac12\coth(\beta_j/2)=\bar n_j+\tfrac12.
\label{eq:mode-temperature}
\end{equation}
 Deriving this is a standard exercise in the Heisenberg algebra together with equation \eqref{eq:Weyl}. One can see that a pure mode has
$\upsilon_j=\frac12$ ($\beta_j=\infty$), while a mixed/thermal one has
$\upsilon_j>\frac12$. The vacuum restricted to a finite canonical system
satisfies the strict uncertainty relation
\begin{equation}
V+\tfrac\ii2J_n>0,\quad\text{equivalently}\quad\upsilon_j>\tfrac12
\text{ for all }j,
\label{eq:strict-uncertainty}
\end{equation}
because $V+\frac\ii2J_n$ is the Gram matrix of the one-particle
vectors of $\sx_1,\dots,\sx_{2n}$, which are linearly independent by
Lemma~\ref{lem:standard} (see \eqref{eq:gram}). Hence $\rho_F$ is
\emph{faithful}: it has no zero eigenvalues.

\begin{lemma}[Passage to all local effects]\label{lem:exhaustion}
Let $h$ be continuous on $[0,1]^3$, and suppose that
\[
h\bigl(\omega(E),\omega(P),\omega(EP)\bigr)\le0
\]
for \emph{every} finite canonical system $F_A$ in $B_1(0)$, every
finite canonical system $F_B$ in $B_1(s\eone)$, every effect $E$ of
the algebra of $F_A$ and every effect $P$ of the algebra of $F_B$.
Then the inequality holds for all effects $E\in\A$, $P\in\B$. The
same holds for a bound on $\abs{\omega(AB)-\omega(A)\omega(B)}$ over
self-adjoint contractions.
\end{lemma}

\begin{proof}
The union of the algebras of all finite canonical systems in $B_1(0)$
is a $*$-algebra
(two systems lie in a third system) containing the generators of $\A$, so it
is strongly dense in $\A$. By Kaplansky's density theorem
\cite[Theorem~2.4.16]{BR87}, every self-adjoint contraction of $\A$ is
a strong limit of self-adjoint contractions of the union; applied to
$2E-I$ this gives effects $E_\alpha\to E$ strongly, each belonging to
the algebra of some finite canonical system, and likewise
$P_\beta\to P$ in $\B$. Products of bounded strongly convergent nets
converge strongly, and the vacuum, a vector state, is strongly
continuous (as is any normal state, on bounded sets), so
$\bigl(\omega(E_\alpha),\omega(P_\beta),\omega(E_\alpha P_\beta)\bigr)
\to\bigl(\omega(E),\omega(P),\omega(EP)\bigr)$, and by the
continuity of $h$ the closed condition $h\le0$ passes to the limit.
\end{proof}

\section{A probability criterion for the classical bound}
\label{sec:probability}

In this section $\omega$ is any state on the algebra generated by two
commuting von Neumann algebras $\A,\B$. For $r>1$, the \emph{event comparison
constants} are
\begin{equation}
K_{r,+}=\sup_{E,P}\frac{\omega(EP)^r}{(\omega(E)\omega(P))^{r-1}},
\quad
K_{r,-}=\sup_{E,P}\frac{(\omega(E)\omega(P))^r}{\omega(EP)^{r-1}},
\label{eq:K}
\end{equation}
over projections $E\in\A$, $P\in\B$ with $\omega(E)\omega(P)>0$; a
positive numerator over a zero denominator counts as $+\infty$. Both
constants are at least $1$ ($E=P=I$). For $r_+>1$ and $r_->1$, taking
roots in \eqref{eq:K} gives, for all projections,
\begin{equation}
\begin{gathered}
\omega(EP)\le K_{r_+,+}^{1/r_+}\,
\bigl(\omega(E)\omega(P)\bigr)^{1-1/r_+},\\
\omega(EP)\ge K_{r_-,-}^{-1/(r_--1)}\,
\bigl(\omega(E)\omega(P)\bigr)^{r_-/(r_--1)};
\end{gathered}
\label{eq:event-orders}
\end{equation}
if $\omega(E)\omega(P)=0$ then $\omega(EP)=0$ since $EP\le E,P$. Thus a
finite $K_{r,+}$ says that a joint event is never much more probable
than a power of the product of its marginals, and a finite $K_{r,-}$
says it is never much less probable.

\begin{theorem}[Two event bounds and clustering]\label{thm:probability}
Let $r_+>2$, $1<r_-<2$, and
\begin{equation}
\begin{gathered}
b=1-\frac1{r_+},\quad a=1-\frac1{r_-},\quad
\eta=1-\frac{1-b}{ab^2},\\
\nu_+=\frac{1+b}{ar_+b^2},\quad \nu_-=\frac{1-a}{a}.
\end{gathered}
\label{eq:probability-exponents}
\end{equation}
Suppose $\eta>0$, $K_{r_+,+},K_{r_-,-}<\infty$, and, for some
$q_{\mathrm{cl}}\ge0$,
\begin{equation}
\abs{\omega(AB)-\omega(A)\omega(B)}\le q_{\mathrm{cl}}
\label{eq:additive-cluster}
\end{equation}
for all self-adjoint contractions $A\in\A$, $B\in\B$. If
\begin{equation}
q_{\mathrm{cl}}^{\eta}K_{r_+,+}^{\nu_+}K_{r_-,-}^{\nu_-}\le2^{-2/(ab)},
\label{eq:probability-cutoff}
\end{equation}
then $\beta(\omega;\A,\B)=1$.
\end{theorem}

The proof is in Appendix~\ref{app:probability}. In outline: a
self-adjoint contraction is an average of reflections and the CHSH
expression is affine in each argument, so a violation by contractions
implies a violation by four reflections $A_0,A_1,B_0,B_1$. Their
spectral projections $E_i^\pm=(I\pm A_i)/2$ and $P_j^\pm=(I\pm B_j)/2$
define sixteen joint probabilities $\omega(E_i^\pm P_j^\pm)$, and a
violation forces the least likely of the eight that enter the CHSH
value with a positive sign---by the symmetries of the expression, this
cell can be taken to be $E_0^-P_0^-$---to satisfy two incompatible
bounds: the event bounds \eqref{eq:event-orders} give
\begin{equation}
\omega(E_0^-P_0^-)^{\eta}\,K_{r_+,+}^{\nu_+}K_{r_-,-}^{\nu_-}
>2^{-2/(ab)},
\label{eq:outline-lower}
\end{equation}
while \eqref{eq:additive-cluster} forces
$\omega(E_0^-P_0^-)<q_{\mathrm{cl}}$, and together these contradict
\eqref{eq:probability-cutoff}.

\emph{What remains.---}Theorem~\ref{thm:probability} reduces
Theorem~\ref{thm:main} to verifying its hypotheses for the vacuum and
the two algebras $\A,\B$ at large separation $s$: bounds on
$K_{r_+,+},K_{r_-,-}$ tending to $1$, and a clustering constant
$q_{\mathrm{cl}}$ tending to $0$, fast enough that
\eqref{eq:probability-cutoff} holds for all $s\ge C_*$. Both come from
comparing the vacuum with the product of its marginals.
The comparison degenerates on modes that are nearly pure in the
vacuum; it is rescued by the monogamy stated in
Section~\ref{sec:statement}---purity of a mode in one ball suppresses
its correlation with the other ball at a faster rate.
Section~\ref{sec:gaussian} bounds the quotients in \eqref{eq:K}, for
the vacuum restricted to finitely many modes, by R\'enyi divergences
between the restricted vacuum and the product of its marginals, and
bounds $q_{\mathrm{cl}}$ by their trace distance;
Lemma~\ref{lem:exhaustion} then removes the restriction to finitely
many modes. The divergences are bounded in turn by covariance
differences (Theorem~\ref{thm:gaussian}), which
Section~\ref{sec:estimates} bounds by negative powers of $s$,
uniformly in the mass. Section~\ref{sec:completion} fixes $r_\pm$ and
assembles \eqref{eq:probability-cutoff} with $q_{\mathrm{cl}}$
proportional to $1/s$.

\section{Gaussian comparison and the weight operator}
\label{sec:gaussian}

This section carries out the first step of the programme at the end of
Section~\ref{sec:probability}: for the vacuum restricted to finitely
many modes in the two balls, the quotients in \eqref{eq:K} are bounded
by R\'enyi divergences between the restricted vacuum $\rho_F$ and the
product $\sigma_F$ of its marginals. Furthermore, $q_{\mathrm{cl}}$ is bounded by their
trace distance. Theorem~\ref{thm:gaussian} then bounds the divergences
by covariance differences as in \eqref{eq:gaussian-bound}.
The section then confronts the main difficulty of the argument,
uniformity over the choice of modes, and resolves it by the monogamy
mechanism of Section~\ref{sec:statement}.


\emph{A divergence with data processing.---}For faithful density
operators $\rho,\sigma$ and $r>1$, let
\begin{equation}
\Qtil_r(\rho\|\sigma)=\Tr\bigl(\sigma^{-\tau/2}\rho\,\sigma^{-\tau/2}\bigr)^r,
\qquad\tau=\frac{r-1}r,
\label{eq:renyi}
\end{equation}
with the value $+\infty$ if the operator inside has infinite $r$-norm.
Its logarithm divided by $r-1$ is the sandwiched R\'enyi divergence
\cite{WWY14,ML13}, a one-parameter family of measures of how
distinguishable two states are; $\Qtil_r=1$ if $\rho=\sigma$ and
$\Qtil_r\ge1$ in general. Its basic property is data processing: it
cannot increase under any physical operation \cite{FL13,Beigi13}. We
use two consequences. Here $\rho(M)=\Tr(\rho M)$, and
$\norm T_p=(\Tr\abs T^p)^{1/p}$ are the Schatten norms; $\norm T_1$,
the \emph{trace norm}, is the sum of the singular values.

\begin{lemma}[Effects and trace distance; cf.~{\cite{FL13,Beigi13,WWY14,TKRVV10}}]\label{lem:renyi-consequences}
Let $\rho,\sigma$ be faithful density operators on one Hilbert space,
with $\Qtil_r(\rho\|\sigma)<\infty$. Then every effect $M$ on that
space satisfies
\begin{equation}
\rho(M)^r\le\Qtil_r(\rho\|\sigma)\,\sigma(M)^{r-1},
\label{eq:effect-renyi}
\end{equation}
and, for $r\ge2$,
\begin{equation}
\norm{\rho-\sigma}_1\le\sqrt{\Qtil_r(\rho\|\sigma)^{1/(r-1)}-1}.
\label{eq:trace-distance}
\end{equation}
\end{lemma}

The proof is in Appendix~\ref{app:gaussian}. A word on where density
matrices live: the restrictions of the vacuum to the local
algebras $\A,\B$ are normal states, but these algebras are von Neumann
algebras of type III, without a trace, so the restricted states have
no intrinsic density matrix there. Density matrices exist on the
finite canonical systems of Section~\ref{sec:field}, which are
ordinary quantum mechanics by the Stone--von Neumann theorem; the
density matrices used below represent the vacuum restricted to such
systems, and every trace below is taken there. In the end, the resulting inequalities will be between
the numbers $\omega(E)$, $\omega(P)$, $\omega(EP)$ with constants
independent of the finite system, and so the limits can be uniformly taken to $\A,\B$, by
Lemma~\ref{lem:exhaustion}.

Returning to the vacuum
$\omega$, let $F_A,F_B$ be
finite canonical systems (finitely many modes) in $B_1(0)$ and $B_1(s\eone)$. Since
$\varsigma$ vanishes between them, $F=F_A\oplus F_B$ is a finite
canonical system whose two halves generate commuting algebras. Let
$\rho_F$ be the vacuum on $F$ and
\[
\sigma_F=\rho_{F_A}\otimes\rho_{F_B},
\]
the product of the vacuum restrictions to the two halves (the
marginals of $\rho_F$), both faithful Gaussian states by
\eqref{eq:strict-uncertainty}. For effects $E$ of $F_A$ and $P$ of
$F_B$, $EP$ is an effect with $\rho_F(EP)=\omega(EP)$ and
$\sigma_F(EP)=\omega(E)\omega(P)$. Applying
\eqref{eq:effect-renyi} to $(\rho_F,\sigma_F,r_+)$ and to
$(\sigma_F,\rho_F,r_-)$ with $M=EP$,
\begin{equation}
\begin{gathered}
\omega(EP)^{r_+}\le\Qtil_{r_+}(\rho_F\|\sigma_F)\,
\bigl(\omega(E)\omega(P)\bigr)^{r_+-1},\\
\bigl(\omega(E)\omega(P)\bigr)^{r_-}\le\Qtil_{r_-}(\sigma_F\|\rho_F)\,
\omega(EP)^{r_--1}.
\end{gathered}
\label{eq:finite-events}
\end{equation}
These are exactly the quotients of \eqref{eq:K}. The task is to bound
$\log\Qtil_{r_+}(\rho_F\|\sigma_F)$ and
$\log\Qtil_{r_-}(\sigma_F\|\rho_F)$ by quantities that vanish as
$s\to\infty$, uniformly in $\mu$ and in $F_A,F_B$.

\emph{Thermal powers.---}If $\rho_V$ is a faithful Gaussian state with
covariance $V$ and $0<\tau<1$, then $\rho_V^\tau/\Tr\rho_V^\tau$ is
Gaussian with covariance $V^{[\tau]}$, which has the same normal modes
as $V$ and the symplectic eigenvalues \cite[Proposition~1]{SLW18}
\begin{equation}
\upsilon_j^{[\tau]}=\tfrac12\coth\Bigl(\tau\artanh\frac1{2\upsilon_j}\Bigr),
\label{eq:thermal-power}
\end{equation}
the $\tau$th power of a thermal state being the thermal state at
$\tau$ times the inverse temperature. The \emph{thermal increment}
$\Therm_\tau(V)=V^{[\tau]}-V$ is positive definite, since $\coth$
decreases.

\begin{theorem}[Gaussian R\'enyi comparison]\label{thm:gaussian}
Let $\rho,\sigma$ be centred (zero-mean) faithful Gaussian states of the same $n$
modes, with covariance matrices $V_\rho,V_\sigma$. For $r>1$ let
\[
\begin{gathered}
\tau=\frac{r-1}r,\qquad k_r=\max\{2,\tau^{-1}\},\\
\mathsf Z=\Therm_\tau(V_\sigma)^{-1/2}(V_\rho-V_\sigma)\Therm_\tau(V_\sigma)^{-1/2}.
\end{gathered}
\]
If $\norm{\mathsf Z}\le(4k_r)^{-1}$, then
\begin{equation}
\log\Qtil_r(\rho\|\sigma)\le2rk_r\norm{\mathsf Z}_1 .
\label{eq:gaussian-bound}
\end{equation}
\end{theorem}

The proof is in Appendix~\ref{app:gaussian}. Neither $n$ nor the
symplectic eigenvalues appear: the covariance difference enters only
relative to the thermal increment of $\sigma$.

\emph{The obstruction: nearly pure modes.---}Theorem~\ref{thm:gaussian}
and \eqref{eq:finite-events} cannot be combined as they stand to prove the main result.
Lemma~\ref{lem:exhaustion} demands constants independent of the chosen
modes $F$, so $\norm{\mathsf Z}_1$ must first be bounded uniformly in
$F$---and $\mathsf Z$ is the covariance difference \emph{divided by}
the thermal increment of $\sigma_F$. Consider a single normal mode of
$\sigma_F$, thermal at inverse temperature $\beta_j$ as in
\eqref{eq:mode-temperature}. By \eqref{eq:thermal-power} and \eqref{eq:mode-temperature}, the thermal increment
on the mode takes the form
\[
\begin{gathered}
\upsilon_j^{[\tau]}-\upsilon_j
=\tfrac12\coth(\tau\beta_j/2)-\tfrac12\coth(\beta_j/2)\\
\sim e^{-\tau\beta_j},\qquad\beta_j\to\infty,
\end{gathered}
\]
and so on that mode each side of $\mathsf Z$ divides the covariance
difference by a quantity of order $e^{-\tau\beta_j/2}$. Nothing stops $F$ from
containing modes with $\beta_j$ as large as one likes (physically,
vacuum entanglement concentrates near the boundary of a region, and
modes supported well inside a ball are nearly pure), so no bound on
the covariance difference alone controls $\norm{\mathsf Z}_1$
uniformly in $F$.

\emph{The rescue: purity is monogamous.---}So what must be shown
is that the covariance difference itself decays on nearly pure modes,
at least as fast as $e^{-\tau\beta_j/2}$. Since
$\sigma_F$ is the product of the marginals of $\rho_F$, the diagonal
blocks of $V_\rho-V_\sigma$ cancel: what remains is the covariance
between the two balls. And the vacuum is a pure state: across the
split between the mode and the environment, the vacuum vector has a
Schmidt decomposition with the thermal weights of
\eqref{eq:mode-temperature},
\[
\Omega=\sum_{n\ge0}\sqrt{p_n}\;|n\rangle\otimes|\widetilde
n\rangle,\qquad p_n=(1-e^{-\beta_j})\,e^{-\beta_jn},
\]
with $|n\rangle$ the $n$-quanta state of the mode and
$|\widetilde n\rangle$ orthonormal vectors of the environment, which
here means everything other than the mode. Let $q$ be the position of
the mode, which connects $|n\rangle$ only to $|n\pm1\rangle$, with
$\langle n+1|q|n\rangle=\sqrt{(n+1)/2}$, and let $\sy$ be any source
in the other ball, so that the smeared field $\Phi(\sy)$ acts on the
environment factor alone, with $\omega(\Phi(\sy))=0$. The covariance
is
\[
\begin{aligned}
\omega(q\,\Phi(\sy))&=\sum_{n}\sqrt{p_np_{n+1}}\,\sqrt{2(n+1)}\,
\operatorname{Re}\langle\widetilde{n+1}|\Phi(\sy)|\widetilde n\rangle\\
&=e^{-\beta_j/2}\sum_np_n\sqrt{2(n+1)}\,
\operatorname{Re}\langle\widetilde{n+1}|\Phi(\sy)|\widetilde n\rangle,
\end{aligned}
\]
because $p_{n+1}=e^{-\beta_j}p_n$: each term pays
$e^{-\beta_j/2}$ for moving one step along the Schmidt sum. The
Cauchy--Schwarz inequality and
$\sum_np_n\langle\widetilde n|\Phi(\sy)^2|\widetilde n\rangle
=\omega(\Phi(\sy)^2)=\cov_\mu(\sy,\sy)$ give
\[
\begin{aligned}
\abs{\omega(q\,\Phi(\sy))}
&\le e^{-\beta_j/2}\sqrt{2(\bar n_j+1)}\;\cov_\mu(\sy,\sy)^{1/2}\\
&=\sqrt{2\bar n_j}\;\cov_\mu(\sy,\sy)^{1/2},
\end{aligned}
\]
where $\bar n_j=\sum_nnp_n=1/(e^{\beta_j}-1)$ is the mean
number of quanta in the mode, as in \eqref{eq:mode-temperature}, so
that $e^{-\beta_j}(\bar n_j+1)=\bar n_j$; likewise for the momentum. A
pure mode ($\bar n_j=0$) is
therefore correlated with nothing---a pure state with a pure marginal
factorizes---and every covariance of a nearly pure mode with its
environment is suppressed by $\sqrt{2\bar n_j}\sim\sqrt2\,e^{-\beta_j/2}$.
Since $\tau<1$, the
covariance difference decays faster than the conjugation by
$\Therm_\tau(V_\sigma)^{-1/2}$ in $\mathsf Z$ grows.

Showing this rigorously is the
plan for the rest of the proof. The operator $\Res$, constructed
next, acts on the full source space of a ball and on mode $j$ as a
multiple of $e^{-\beta_j/2}$: it is the per-mode factor of the
covariance bound above, made independent of $F$. Dividing this factor
out of the covariance difference on each side defines, schematically,
\[
Y=\Res^{-\tau}(V_\rho-V_\sigma)\Res^{-\tau},\quad\text{i.e.}\quad
V_\rho-V_\sigma=\Res^{\,\tau}Y\Res^{\,\tau}.
\]
Since $\Res$ has eigenvalues accumulating at $0$, $\Res^{-\tau}$
is unbounded and $Y$ could a priori be infinite. Section~\ref{sec:estimates}
(Theorems \ref{thm:forward} and~\ref{thm:reverse}) shows that it is
not: $Y$ is trace class, with $\norm Y_1$ decaying in $s$ and
independent of $F$---the surplus $e^{-(1-\tau)\beta_j/2}$ per mode is
what pays for this. In $\mathsf Z$,
each factor $\Res^{\,\tau}$ then cancels one factor
$\Therm_\tau(V_\sigma)^{-1/2}$: Lemma~\ref{lem:thermal-comparison}
below shows that the two vanish at the same rate,
Lemma~\ref{lem:compression} shows that restricting to $F$ does not
spoil this, and the resulting bounds \eqref{eq:one-space} and
\eqref{eq:trace-compression} on $\norm{\mathsf Z}_1$ do not depend on
$F$.

\emph{The weight operator.---}We now construct $\Res$, first on
the $n$ modes of $F$. Choose coordinates in which $V_\sigma=I$, and
let $\mathsf J$ be the real antisymmetric matrix of commutators in
these coordinates, $[\Phi(\sx_k),\Phi(\sx_l)]=\ii\mathsf J_{kl}$,
and let
\[
G=I+\tfrac\ii2\mathsf J,\qquad\overline G=I-\tfrac\ii2\mathsf J.
\]
In both applications below, $\sigma$ is the vacuum restricted to
the modes in question (to the modes of one ball, or to all of $F$),
so $V_\sigma=I$ means $\cov_\mu(\sx_k,\sx_l)=\delta_{kl}$, and by
\eqref{eq:gram} $G_{kl}=\ip{\kappa\sx_k}{\kappa\sx_l}$ is the matrix
of inner products of the one-particle vectors. In
normal-mode coordinates rescaled to $V_\sigma=I$, mode $j$ has
covariance $I_2$ and commutator block $J_1/\upsilon_j$, so $G$ and
$\overline G$ have eigenvalues $1\pm d_j$, with
\[
d_j=\frac1{2\upsilon_j}=\tanh(\beta_j/2)\in(0,1)
\]
by \eqref{eq:mode-temperature} and \eqref{eq:strict-uncertainty}.
Write $d$ for the matrix $\abs{\frac\ii2\mathsf J}$, with eigenvalues
$d_j$. The \emph{weight operator} is
\begin{equation}
\Res=(G\overline G)^{1/2}=(I-d^2)^{1/2},
\label{eq:reserve-finite}
\end{equation}
which acts on mode $j$ as
$\operatorname{sech}(\beta_j/2)\sim2e^{-\beta_j/2}$: it vanishes
exactly on pure modes, at the rate of the covariance bound above, and
tends to $1$ on strongly mixed ones. In the same coordinates the
thermal increment \eqref{eq:thermal-power} is
\begin{equation}
\Therm_\tau(I)=d\coth(\tau\artanh d)-I ,
\label{eq:thermal-d}
\end{equation}
which acts on mode $j$ as
$\tanh(\beta_j/2)\coth(\tau\beta_j/2)-1\sim2e^{-\tau\beta_j}$, the
same rate as $\Res^{2\tau}$.

Nothing in this construction needs $F$. On the full source space
$\Hh_{\mu,U}$ of an open set $U$, with inner product $\cov_\mu$,
define $\mathsf J$ by $\cov_\mu(\sx,\mathsf J\sy)=\varsigma(\sx,\sy)$;
then, again by \eqref{eq:gram},
\begin{equation}
\begin{gathered}
G_U=\kappa_U^*\kappa_U=I+\tfrac\ii2\mathsf J,\\
\overline G_U=2I-G_U,\qquad
\Res_U=(G_U\overline G_U)^{1/2}.
\end{gathered}
\label{eq:reserve}
\end{equation}
By Lemma~\ref{lem:standard}, $G_U$ and $\overline G_U$ are
injective with spectrum in $[0,2]$, so $0<\Res_U\le I$; now the
spectrum of $\Res_U$ may accumulate at $0$. A finite canonical system
$F\subset U$ sits inside $\Hh_{\mu,U}$ through its inclusion $\iota$,
and the matrix of an operator $X$ on $\Hh_{\mu,U}$ between vectors of
$F$ is $\iota^*X\iota$. The next two lemmas, proved in
Appendix~\ref{app:compression}, compare the thermal increment with
$\Res$, and compare $\Res$ for $F$ with $\iota^*\Res_U\iota$.

\begin{lemma}[Thermal comparison]\label{lem:thermal-comparison}
On any source space, finite or full, and for $0<\tau<1$,
\begin{equation}
\Therm_\tau(I)\ge c_\tau\Res^{2\tau},\qquad
c_\tau=\begin{cases}2^{1-2\tau},&0<\tau\le\frac12,\\
(1-\tau)/\tau,&\frac12\le\tau<1,\end{cases}
\label{eq:scalar-reserve}
\end{equation}
where on the full source space $\Therm_\tau(I)$ is defined by
\eqref{eq:thermal-d} in the functional calculus of $d$.
\end{lemma}

\begin{lemma}[Compression]\label{lem:compression}
Let $\iota$ be the inclusion of a finite canonical system $F$ into a
full source space---isometric for the covariance norm and preserving
$\varsigma$, so that $\iota^*\mathsf J\iota$ is the commutator form of
$F$---and let
$G_F=\iota^*G\iota$, $\overline G_F=\iota^*\overline G\iota$,
$\Res_F=(G_F\overline G_F)^{1/2}$ be the operators \eqref{eq:reserve}
of $F$ itself. For $0<a\le\frac12$,
\begin{equation}
\Res_F^{2a}\ge\iota^*\Res^{2a}\iota ,
\label{eq:low-compression}
\end{equation}
and for an integer $r\ge2$ and $b=(r-1)/r$,
\begin{equation}
\Therm_b(I_F)\ge\iota^*\Therm_b(I)\iota\ge
(r-1)^{-1}\,\iota^*\Res^{2b}\iota .
\label{eq:high-compression}
\end{equation}
\end{lemma}

Together, the two lemmas give the cancellation announced above.
For $0<a\le\frac12$, \eqref{eq:scalar-reserve} on $F$ and
\eqref{eq:low-compression} give
$\Therm_a(I_F)\ge2^{1-2a}\,\iota^*\Res^{2a}\iota$, that is,
$\norm{\Res^a\iota\,\Therm_a(I_F)^{-1/2}}^2\le2^{2a-1}$. Since
$\norm{ABC}_1\le\norm A\norm B_1\norm C$, every trace-class $Y$
satisfies
\begin{equation}
\begin{split}
\bigl\|\Therm_a(I_F)^{-1/2}\iota^*&\Res^aY\Res^a\iota\,
\Therm_a(I_F)^{-1/2}\bigr\|_1\\
&\le2^{2a-1}\norm Y_1 .
\end{split}
\label{eq:one-space}
\end{equation}
In the same way, for the two balls, with weight operators
$\Res_A=\Res_{B_1(0)}$ and $\Res_B=\Res_{B_1(s\eone)}$ and inclusions
$\iota_A\colon F_A\to\mathcal X_{\mu,B_1(0)}$ and
$\iota_B\colon F_B\to\mathcal X_{\mu,B_1(s\eone)}$,
\eqref{eq:high-compression} gives, for an integer $r\ge2$,
$b=(r-1)/r$ and every trace-class $X$,
\begin{equation}
\begin{split}
\bigl\|\Therm_b(I_{F_A})^{-1/2}\iota_A^*&\Res_A^{\,b}X\Res_B^{\,b}
\iota_B\,\Therm_b(I_{F_B})^{-1/2}\bigr\|_1\\
&\le(r-1)\norm X_1 .
\end{split}
\label{eq:trace-compression}
\end{equation}

\section{The two field estimates}\label{sec:estimates}

This section proves the factorization
$V_\rho-V_\sigma=\Res^{\,\tau}Y\Res^{\,\tau}$ of
Section~\ref{sec:gaussian}, with $\norm Y_1$ bounded by a negative
power of $s$ uniformly in $\mu$. It is needed twice, because
\eqref{eq:finite-events} uses two divergences: the \emph{forward}
comparison $\Qtil_{r_+}(\rho_F\|\sigma_F)$, which conjugates by the
thermal increment of $\sigma_F$, and the \emph{reverse} comparison
$\Qtil_{r_-}(\sigma_F\|\rho_F)$, which conjugates by that of $\rho_F$.
For each there is one operator on full source spaces whose matrix on
$F$ is the covariance difference.

For the forward comparison it is the \emph{cross-covariance
operator} $\mathsf C_{\mu,s}\colon\Hh_{\mu,B_1(s\eone)}\to\Hh_{\mu,B_1(0)}$,
defined by
\begin{equation}
\ip\sx{\mathsf C_{\mu,s}\sy}_{\mu,B_1(0)}=\cov_\mu(\sx,\sy).
\label{eq:C-def}
\end{equation}
For the reverse comparison, let $U_s=B_1(0)\cup B_1(s\eone)$ and
split each source on $U_s$ as $\sx=\sx_A+\sx_B$ into its parts in the
two balls. The \emph{product-minus-joint operator} $\JointD_{\mu,s}$
on $\Hh_{\mu,U_s}$, which carries the joint covariance $\cov_\mu$ as
inner product, is defined by deleting the cross terms:
\begin{equation}
\begin{split}
\ip\sx{\JointD_{\mu,s}\sy}_{\mu,U_s}
&=\cov_\mu(\sx_A,\sy_A)+\cov_\mu(\sx_B,\sy_B)-\cov_\mu(\sx,\sy)\\
&=-\cov_\mu(\sx_A,\sy_B)-\cov_\mu(\sx_B,\sy_A).
\end{split}
\label{eq:D-joint}
\end{equation}
Section~\ref{sec:completion} checks that
$\iota_A^*\mathsf C_{\mu,s}\iota_B$ is the off-diagonal block of
$V_{\rho_F}-V_{\sigma_F}$ in coordinates with $V_{\sigma_F}=I$, and
that $\iota_F^*\JointD_{\mu,s}\iota_F=V_{\sigma_F}-V_{\rho_F}$ in
coordinates with $V_{\rho_F}=I$, where $\iota_F$ is the inclusion of
$F$ into $\Hh_{\mu,U_s}$. The trace-class factor called $Y$ in
Section~\ref{sec:gaussian} is $X_{b,\mu,s}$ below for the forward
comparison and $Y_{\mu,s}$ for the reverse one.

\begin{theorem}[Forward estimate]\label{thm:forward}
For each $\frac12\le b<1$ there are finite $S_b,F_b$, independent of
$\mu$, such that for all $\mu>0$ and $s\ge S_b$ there is a trace-class
$X_{b,\mu,s}$ with
\begin{equation}
\mathsf C_{\mu,s}=\Res_A^{\,b}X_{b,\mu,s}\Res_B^{\,b},\qquad
\norm{X_{b,\mu,s}}_1\le F_bs^{-2}.
\label{eq:forward-trace}
\end{equation}
\end{theorem}

\begin{theorem}[Reverse estimate]\label{thm:reverse}
Let $a=2^{-18}$, $\gamma=2(1-2^{-17})$, and let $\JointRes$ be the
weight operator \eqref{eq:reserve} of $U_s$. There are finite
$S_{\mathrm{rev}},\mathcal J_0$, independent of $\mu$, such that for
all $\mu>0$ and $s\ge S_{\mathrm{rev}}$ there is a self-adjoint
trace-class $Y_{\mu,s}$ with
\begin{equation}
\JointD_{\mu,s}=\JointRes^{\,a}Y_{\mu,s}\JointRes^{\,a},
\quad
\norm{Y_{\mu,s}}_1\le\mathcal J_0s^{-\gamma}.
\label{eq:reverse-trace}
\end{equation}
\end{theorem}

As in Section~\ref{sec:gaussian}, the content of each theorem is
that the middle factor exists and is trace class, although
$\Res_A^{-b}$, $\Res_B^{-b}$ and $\JointRes^{-a}$ are unbounded. The
decay is uniform in $\mu$, including $\mu s\to0$; since at $\mu=0$
vacuum correlations decay only as a power of the distance, it can only
be polynomial. The proofs are in Appendix~\ref{app:estimates}. They
rest on the following bound, which controls the one-particle vector
$\kappa\sx$ of a source in a ball, continued to complex momenta, by
the weighted norm $\norm{\Res_U^b\sx}$.

\begin{theorem}[Weighted complex-shell evaluation]\label{thm:evaluation}
For a source $\sx=(f,g)$ let $\Psi_\sx(z)=\hat f(\bz)+\ii z_0\hat g(\bz)$,
$z\in\C^4$, an entire function that equals $\kappa\sx$ on $H_\mu^+$,
and let $H_\mu^{\C}=\{z\in\C^4:z\cdot z=-\mu^2\}$ be the complex
mass shell, where $z\cdot z=-z_0^2+\bz\cdot\bz$ is the Minkowski product
extended bilinearly, without complex conjugation. For every $0<c<\frac12$ there are finite $C_c,\lambda_c\ge1$
such that for all $\mu,T>0$, nonempty bounded open $U\subset B_T(0)$,
$\sx\in\Hh_{\mu,U}$ and every energy scale $\Energy\ge\mu$,
\begin{equation}
\abs{\Psi_\sx(z)}^2\le C_c\Energy^{-2}e^{2\lambda_cT\Energy}
\norm{G_U^{1/2-c}\overline G_U^{\,c}\sx}^2
\label{eq:evaluation}
\end{equation}
for all $z\in H_\mu^{\C}$ with $\abs z_{\mathrm E}\le2\Energy$
(Euclidean norm on $\C^4$); $\lambda_c=2048/\sin(2\pi c)$ is admissible.
\end{theorem}

The right-hand side is a weighted norm of $\sx$. Indeed, let
$\frac12\le b<1$. Since $\Res_U^b=G_U^{b/2}\overline G_U^{\,b/2}$
and $G_U+\overline G_U=2I$,
\[
G_U^{1/2-c}\overline G_U^{\,c}=\begin{cases}
\Res_U^b\,\overline G_U^{\,(1-2b)/2},&c=\frac{1-b}2,\\[2pt]
\Res_U^b\,G_U^{(1-2b)/2},&c=\frac b2,
\end{cases}
\]
and the second factor is at most $1$ on the part of the spectrum
where $G_U\le1$ in the first case and where $G_U\ge1$ in the second.
Using each choice on its part, for the same $z$,
\begin{equation}
\abs{\Psi_\sx(z)}\le\bigl(C_{(1-b)/2}^{1/2}+C_{b/2}^{1/2}\bigr)
\Energy^{-1}e^{\lambda_{b/2}T\Energy}\norm{\Res_U^{\,b}\sx},
\label{eq:symmetric-evaluation}
\end{equation}
with $\lambda_{(1-b)/2}=\lambda_{b/2}=2048/\sin(\pi b)$.

The bound \eqref{eq:evaluation} comes from the
Bisognano--Wichmann theorem, which needs the following definitions.
A closed real
subspace $\mathcal K$ of a complex Hilbert space is \emph{standard} if
$\mathcal K\cap\ii\mathcal K=\{0\}$ and $\mathcal K+\ii\mathcal K$ is
dense; the antilinear map $\mathsf s_{\mathcal K}(u+\ii v)=u-\ii v$,
$u,v\in\mathcal K$, is then well defined and closable, and
$\Delta_{\mathcal K}=\mathsf s_{\mathcal K}^*\mathsf s_{\mathcal K}$
is a positive operator, the \emph{modular operator} of $\mathcal K$
\cite{RvD77}. For a region, take the closed real span of the
one-particle vectors of the real sources it contains: $\mathcal K_U$
for the ball $U$, standard by Lemma~\ref{lem:standard}, and
$\mathcal K_{\mathcal W_\ell}$ for the wedge
$\mathcal W_\ell=\{x^1<\ell-\abs{x^0}\}$, with sources in the base
$\{x^1<\ell\}$. The wedge subspace $\mathcal K_{\mathcal W_\ell}$
is standard too: it contains the ball's
$\mathcal K_U$, so
$\mathcal K_{\mathcal W_\ell}+\ii\mathcal K_{\mathcal W_\ell}$ is
dense; and if $\xi$ and $-\ii\xi$ both lie in
$\mathcal K_{\mathcal W_\ell}$, then for every source $\sy$ in a ball
$U'$ in the opposite half-space $\{x^1>\ell\}$ the numbers
$\ip\xi{\kappa_{U'}\sy}$ and $\ii\ip\xi{\kappa_{U'}\sy}$ are both
real, because $\varsigma$ vanishes between the half-spaces
\eqref{eq:gram}; so $\xi\perp\kappa_{U'}\sy$, these vectors are dense
by Lemma~\ref{lem:standard}, and $\xi=0$.

The Bisognano--Wichmann theorem \cite{BW75}, in its
one-particle form \cite{BGL02}, says that
$\Delta_{\mathcal W_\ell}^{\ii y}$ is the boost of rapidity $2\pi y$
that fixes the edge $\{x^0=0,\,x^1=\ell\}$ of the wedge, that is, the
boost along the first axis conjugated by the translation by
$\ell\eone$. Continued to imaginary $y$,
$\Delta_{\mathcal W_\ell}^{\,c}$ for $0<c<\frac12$ is a boost by the
imaginary rapidity $\ii\varphi$, $\varphi=2\pi c$, and acts on
$\kappa_U\sx$ by evaluating $\Psi_\sx$ at complex momenta; the
translation contributes the factor
$\abs{e^{\ii\ell(p_1\cos\varphi+\ii p_0\sin\varphi)}}^2
=e^{-2\ell p_0\sin\varphi}$:
\begin{equation}
\norm{\Delta_{\mathcal W_\ell}^{\,c}\kappa_U\sx}^2
=\int_{H_\mu^+}e^{-2\ell p_0\sin\varphi}
\abs{\Psi_\sx(b_\varphi(p))}^2\dd\mathfrak m_\mu(p),
\label{eq:wedge-contour}
\end{equation}
with
\[
b_\varphi(p)=(p_0\cos\varphi+\ii p_1\sin\varphi,\ p_1\cos\varphi+\ii p_0\sin\varphi,\ p_2,\ p_3).
\]
To bring in $\Res_U$, let $U\subset B_T(0)$ with $T\le\ell$, so
that $D_T\subset\mathcal W_\ell$ and
$\mathcal K_U\subset\mathcal K_{\mathcal W_\ell}$. Then
$\mathsf s_{\mathcal K_U}$ is a restriction of
$\mathsf s_{\mathcal K_{\mathcal W_\ell}}$, so
$\norm{\Delta_{\mathcal W_\ell}^{1/2}\xi}=\norm{\Delta_{\mathcal K_U}^{1/2}\xi}$
on the domain of $\mathsf s_{\mathcal K_U}$, and
$\Delta_{\mathcal W_\ell}\le\Delta_{\mathcal K_U}$ as quadratic forms.
Because $t\mapsto t^{2c}$ is operator monotone for $0<2c<1$
(L\"owner--Heinz \cite{Heinz51}),
$\Delta_{\mathcal W_\ell}^{2c}\le\Delta_{\mathcal K_U}^{2c}$. Finally,
with the polar decomposition $\kappa_U=\mathcal U_UG_U^{1/2}$, one has
$\Delta_{\mathcal K_U}=\mathcal U_U\overline G_UG_U^{-1}\mathcal U_U^*$
(Appendix~\ref{app:evaluation}), so
\begin{equation}
\norm{\Delta_{\mathcal W_\ell}^{\,c}\kappa_U\sx}
\le\norm{G_U^{1/2-c}\overline G_U^{\,c}\sx}.
\label{eq:inclusion}
\end{equation}
Together, \eqref{eq:wedge-contour} and \eqref{eq:inclusion} bound
an integral of $\abs{\Psi_\sx}^2$ over the complex momenta $b_\varphi(p)$,
$p\in H_\mu^+$, by $\norm{G_U^{1/2-c}\overline G_U^{\,c}\sx}^2$.
Appendix~\ref{app:evaluation} turns this into the pointwise bound
\eqref{eq:evaluation} in three steps. (i) Repeating the bound for a
compact family of Lorentz transforms of the wedge and averaging bounds
the integral of $\abs{\Psi_\sx}^2$ over an open patch of
$H_\mu^{\C}$. (ii) $\Psi_\sx$ is holomorphic, so $\abs{\Psi_\sx}^2$ at
a point is at most its average over a small neighbourhood in
$H_\mu^{\C}$; this bounds $\abs{\Psi_\sx}^2$ on the patch. (iii) For
the same reason, $\abs{\Psi_\sx}^2$ at the centre of a complex disc
lying in $H_\mu^{\C}$ is at most its average over the boundary circle;
every $z\in H_\mu^{\C}$ with $\abs z_{\mathrm E}\le2\Energy$ is the
centre of such a disc whose boundary lies in the patch.

\emph{Outline of the proofs of Theorems \ref{thm:forward}
and~\ref{thm:reverse}.---}At the two shell points
$z_\pm=(\pm\sqrt{\mu^2+\bz\cdot\bz},\bz)$ one has
$\hat f(\bz)=\frac12({\Psi_\sx(z_+)+\Psi_\sx(z_-)})$, and similarly for
$\hat g$, so \eqref{eq:symmetric-evaluation} bounds $\hat f$ and
$\hat g$ on $\C^3$. Since the derivatives of $\hat f$ at $0$ are the
moments of $f$, Cauchy estimates then bound the moments of a source
about a centre $\bx_0$,
\[
\Bigl|\int(\bx-\bx_0)^\alpha f\Bigr|
\le M_b(B_bT)^{\abs\alpha}e^{\Lambda_b\mu T}\,\norm{\Res_U^{\,b}\sx},
\]
and similarly for $g$, for every multi-index $\alpha\in\N^3$ and
$U\subset B_T(\bx_0)$, with constants $M_b,B_b,\Lambda_b$ depending
only on $b$ (Lemma~\ref{lem:moments}, which also keeps a factor
$T/(1+\mu T)$ omitted here). Between the two balls the covariance is an
integral against a smooth kernel, a Bessel function of $\by-\bx$
(Appendix~\ref{app:estimates}). For the forward estimate, expand this
kernel in a Taylor series in $\bx$ and $\by$ about the two centres:
\[
\cov_\mu(\sx,\sy)=\sum_{\alpha,\beta}c_{\alpha\beta}
\int\bx^\alpha f\int\by^\beta f'+(\text{same for }g,g').
\]
Each moment of $\sx$ is bounded by $\norm{\Res_A^b\sx}$ and each
moment of $\sy$ by $\norm{\Res_B^b\sy}$, so each term is a rank-one
operator of the form $\Res_A^b(\cdot)\Res_B^b$, and the sum $X_{b,\mu,s}$
of the middle factors converges in trace norm. For the reverse
estimate, $\JointD_{\mu,s}$ has the kernel $-\cov_\mu$ between the two
balls and $0$ within each ball (by \eqref{eq:D-joint}), so it is not a
single Taylor series. It is approximated by finite-rank operators
$\mathsf F_n$ with polynomial kernels, and each increment
$\mathsf G_n=\mathsf F_{n+1}-\mathsf F_n$ satisfies two bounds: a
decaying one, with $\norm{\mathsf G_n}_1$ of order
$s^{-2}(c_{\mathrm{pol}}/s)^n$ for a constant $c_{\mathrm{pol}}$, up
to factors polynomial in $n$ and bounded in $\mu$, and a
weighted one,
$\mathsf G_n=\JointRes^{1/2}\mathsf B_n\JointRes^{1/2}$ with
$\norm{\mathsf B_n}_1$ growing geometrically in $n$. Interpolating
between them (Hadamard's three-lines theorem) gives, for
$0<\theta<1$,
\[
\norm{\JointRes^{-\theta/2}\mathsf G_n\JointRes^{-\theta/2}}_1
\le\norm{\mathsf G_n}_1^{1-\theta}\norm{\mathsf B_n}_1^{\theta}.
\]
The growth of $\norm{\mathsf B_n}_1$ must be overcome by the decay
of $\norm{\mathsf G_n}_1$, which forces $\theta$ to be small; the
proof takes $\theta=2^{-17}$, which gives the power $a=\theta/2=2^{-18}$
and the decay $s^{-2(1-\theta)}=s^{-\gamma}$.

\section{Completion of the proof}\label{sec:completion}

This section fixes the R\'enyi orders $r_\pm$; combines
Theorems \ref{thm:forward} and~\ref{thm:reverse} with
Section~\ref{sec:gaussian} to bound $K_{r_\pm,\pm}$
\eqref{eq:full-K} and the clustering constant, $q_{\mathrm{cl}}=Q_0/s$
\eqref{eq:cluster}; checks the cutoff \eqref{eq:probability-cutoff} of
Theorem~\ref{thm:probability} for $s\ge C_*$; and so proves
Theorem~\ref{thm:main} and Corollary~\ref{cor:fixed}.

\emph{Parameters.---}With $a=2^{-18}$ from Theorem~\ref{thm:reverse},
fix
\begin{equation}
\begin{gathered}
r_-=\frac1{1-a}=\frac{2^{18}}{2^{18}-1},\qquad
r_+=\frac4a=2^{20},\\
b=1-\frac1{r_+}=1-2^{-20}.
\end{gathered}
\label{eq:orders}
\end{equation}
Then $1<r_-<2$, $r_+$ is an integer as \eqref{eq:high-compression}
requires, $a=1-1/r_-$ and $b=1-1/r_+$ as in
\eqref{eq:probability-exponents}, and
\begin{equation}
\begin{gathered}
\eta=1-\frac{1-b}{ab^2}=1-\frac1{4b^2}>0,\\
\nu_+=\frac{1+b}{4b^2},\qquad \nu_-=2^{18}-1 .
\end{gathered}
\label{eq:completion-exponents}
\end{equation}
The small reverse exponent forces the large forward order: $\eta>0$
needs $r_+>1/(ab^2)$. With $F_b$ from Theorem~\ref{thm:forward} and
$\mathcal J_0$ from Theorem~\ref{thm:reverse}, put
\begin{equation}
\varepsilon_+(s)=(r_+-1)F_b\,s^{-2},\qquad
\varepsilon_-(s)=2^{2a-1}\mathcal J_0\,s^{-\gamma}.
\label{eq:envelopes}
\end{equation}

\emph{Forward comparison.---}Choose coordinates on
$F=F_A\oplus F_B$ in which $V_{\sigma_F}=I$; since
$\sigma_F=\rho_{F_A}\otimes\rho_{F_B}$, these are coordinates on
$F_A$ and on $F_B$ orthonormal for $\cov_\mu$, and the inclusions
$\iota_A\colon F_A\to\mathcal X_{\mu,B_1(0)}$ and
$\iota_B\colon F_B\to\mathcal X_{\mu,B_1(s\eone)}$ are isometric. By
\eqref{eq:C-def},
\[
V_{\rho_F}-V_{\sigma_F}=\begin{pmatrix}0&\iota_A^*\mathsf C_{\mu,s}\iota_B\\
(\iota_A^*\mathsf C_{\mu,s}\iota_B)^*&0\end{pmatrix}.
\]
The thermal increment of the product $\sigma_F$ at $\tau=b$ is
block diagonal, with blocks $\Therm_b(I_{F_A})$ and
$\Therm_b(I_{F_B})$, so the matrix $\mathsf Z$ of
Theorem~\ref{thm:gaussian} for $(\rho_F,\sigma_F,r_+)$, which we call
$\mathsf Z_{+,F}$, has zero diagonal blocks and the off-diagonal
block
\[
\mathsf Z_{AB}=\Therm_b(I_{F_A})^{-1/2}\,
\iota_A^*\mathsf C_{\mu,s}\iota_B\,\Therm_b(I_{F_B})^{-1/2};
\]
its norms do not depend on the choice of coordinates
(Appendix~\ref{app:gaussian}). For $s\ge S_b$,
\eqref{eq:forward-trace} and \eqref{eq:trace-compression} with
$r=r_+$ give
$\norm{\mathsf Z_{AB}}_1\le(r_+-1)F_bs^{-2}=\varepsilon_+(s)$,
and since the eigenvalues of $\mathsf Z_{+,F}$ are $\pm$ the singular
values of $\mathsf Z_{AB}$ and zeros,
\begin{equation}
\norm{\mathsf Z_{+,F}}\le\varepsilon_+(s),\qquad
\norm{\mathsf Z_{+,F}}_1\le2\varepsilon_+(s).
\label{eq:finite-forward}
\end{equation}
Here $k_{r_+}=\max\{2,1/b\}=2$, so if $\varepsilon_+(s)\le\frac18$,
Theorem~\ref{thm:gaussian} gives
\begin{equation}
\log\Qtil_{r_+}(\rho_F\|\sigma_F)\le8r_+\varepsilon_+(s).
\label{eq:Qforward}
\end{equation}

\emph{Reverse comparison.---}Now choose coordinates on $F$ in
which $V_{\rho_F}=I$, that is, orthonormal for the joint covariance
$\cov_\mu$, so that the inclusion $\iota_F\colon F\to\mathcal X_{\mu,U_s}$
is isometric. By \eqref{eq:D-joint},
$V_{\sigma_F}-V_{\rho_F}=\iota_F^*\JointD_{\mu,s}\iota_F$, so the
matrix $\mathsf Z$ of Theorem~\ref{thm:gaussian} for
$(\sigma_F,\rho_F,r_-)$ is
\[
\mathsf Z_{-,F}=\Therm_a(I_F)^{-1/2}\,\iota_F^*\JointD_{\mu,s}\iota_F\,\Therm_a(I_F)^{-1/2},
\]
and for $s\ge S_{\mathrm{rev}}$, \eqref{eq:reverse-trace} and
\eqref{eq:one-space} with $\Res=\JointRes$ give
\begin{equation}
\norm{\mathsf Z_{-,F}}_1\le2^{2a-1}\mathcal J_0s^{-\gamma}=\varepsilon_-(s).
\label{eq:finite-reverse}
\end{equation}
Here $k_{r_-}=1/a$, so if $\varepsilon_-(s)\le a/4$,
\begin{equation}
\log\Qtil_{r_-}(\sigma_F\|\rho_F)\le\frac{2r_-}a\varepsilon_-(s).
\label{eq:Qreverse}
\end{equation}
Neither bound depends on $F_A$, $F_B$ or the temperatures
$\beta_j$ of their modes.

\emph{All local events.---}Inserting \eqref{eq:Qforward} and
\eqref{eq:Qreverse} into \eqref{eq:finite-events},
\begin{equation}
\begin{gathered}
\omega(EP)^{r_+}\le e^{8r_+\varepsilon_+(s)}
\bigl(\omega(E)\omega(P)\bigr)^{r_+-1},\\
\bigl(\omega(E)\omega(P)\bigr)^{r_-}\le e^{2r_-\varepsilon_-(s)/a}\,
\omega(EP)^{r_--1},
\end{gathered}
\label{eq:events}
\end{equation}
for effects of the finite algebras. Both are free of denominators and
continuous in the three probabilities, with constants depending only
on $s$, so Lemma~\ref{lem:exhaustion} extends them to all $E\in\A$,
$P\in\B$; for projections with $\omega(E)\omega(P)>0$ they bound the
quotients in \eqref{eq:K}, the second forcing $\omega(EP)>0$. Hence
the event comparison constants satisfy
\begin{equation}
\log K_{r_+,+}\le8r_+\varepsilon_+(s),\qquad
\log K_{r_-,-}\le\frac{2r_-}a\varepsilon_-(s).
\label{eq:full-K}
\end{equation}
For the correlations, \eqref{eq:trace-distance} with $r_+\ge2$ and
\eqref{eq:Qforward} give
\[
\norm{\rho_F-\sigma_F}_1\le\sqrt{\exp\Bigl(\frac{8\varepsilon_+(s)}b\Bigr)-1},
\]
using $r_+/(r_+-1)=1/b$; for contractions $A,B$ of the finite
algebras, $\omega(AB)-\omega(A)\omega(B)=\Tr((\rho_F-\sigma_F)AB)$ is
bounded by this, and Lemma~\ref{lem:exhaustion} extends the bound to
$\A,\B$. If $\varepsilon_+(s)\le b/8$, then $e^x-1\le2x$ on $[0,1]$
gives
\begin{equation}
\abs{\omega(AB)-\omega(A)\omega(B)}\le\frac{Q_0}s,\qquad
Q_0=4\sqrt{(r_+-1)F_b/b}.
\label{eq:cluster}
\end{equation}

\emph{The constant.---}Define
\begin{align}
C_{\mathrm{prob}}&=Q_0\exp\Bigl(\frac{2\log2/(ab)+\nu_++\nu_-}{\eta}\Bigr),
\label{eq:Cprob}\\
C_*&=\max\Bigl\{5,\ S_b,\ S_{\mathrm{rev}},\ \sqrt{8r_+(r_+-1)F_b},
\nonumber\\
&\qquad\qquad
\Bigl(\frac{2^{2a+1}\mathcal J_0}{a}\Bigr)^{1/\gamma},\ C_{\mathrm{prob}}\Bigr\}.
\label{eq:Cstar}
\end{align}
Every entry is finite and independent of $\mu$. For the parameters
\eqref{eq:orders}, $C_{\mathrm{prob}}=Q_0\,10^{362\,232.284166\ldots}$;
the power of ten is an explicit multiplier relative to $Q_0$, not
a numerical determination of $C_*$:
$F_b$, $\mathcal J_0$ and hence $Q_0$ are defined through finite
covers and approximation bounds and are known only to be finite. The
first entry gives $C_*\ge5>2$ independently of $Q_0$, and the proof of
Theorem~\ref{thm:forward} takes
$S_b\ge(4+80\sqrt3)\,\Lambda_b$ with
$\Lambda_b=2^{15}/\sin(\pi b)$ as in \eqref{eq:Lambda}
(Appendix~\ref{app:estimates}), which at $b=1-2^{-20}$ exceeds
$1.5\times10^{12}$.

\begin{proof}[Proof of Theorem~\ref{thm:main}]
Fix $\mu>0$ and $s\ge C_*$. Then $s\ge5>2$, $s\ge S_b$ and
$s\ge S_{\mathrm{rev}}$, so \eqref{eq:finite-forward} and
\eqref{eq:finite-reverse} hold for all finite systems. From
$s^2\ge8r_+(r_+-1)F_b$, $\varepsilon_+(s)\le1/(8r_+)\le b/8<\frac18$,
using $r_+b=r_+-1\ge1$; from $s^\gamma\ge2^{2a+1}\mathcal J_0/a$,
$\varepsilon_-(s)\le a/4$.
Hence \eqref{eq:Qforward}, \eqref{eq:Qreverse}, \eqref{eq:full-K} and
\eqref{eq:cluster} hold, with $\log K_{r_+,+}\le1$ and
$\log K_{r_-,-}\le r_-/2<1$. Take $q_{\mathrm{cl}}=Q_0/s$; since
$s\ge C_{\mathrm{prob}}$,
\[
\Bigl(\frac{Q_0}s\Bigr)^\eta K_{r_+,+}^{\nu_+}K_{r_-,-}^{\nu_-}
\le\Bigl(\frac{Q_0}{C_{\mathrm{prob}}}\Bigr)^\eta e^{\nu_++\nu_-}
=2^{-2/(ab)},
\]
which is \eqref{eq:probability-cutoff}. The algebras $\A,\B$ commute,
$r_+>2$, $1<r_-<2$ and $\eta>0$, so Theorem~\ref{thm:probability} gives
$\beta_\mu(1,s)=1$. For $m,R>0$ and $L\ge C_*R$,
Lemma~\ref{lem:scaling} gives $\beta_m(R,L)=\beta_{mR}(1,L/R)=1$, and
$C_*$ does not depend on $m,R$.
\end{proof}

\begin{proof}[Proof of Corollary~\ref{cor:fixed}]
Let $e$ be the unit vector of the given direction and $\mathcal L$ a
Lorentz transformation with $\mathcal Le=(0,\eone)$. The vacuum is
Poincar\'e invariant, and the unitary implementing $\mathcal L$ carries
$\A_m(O)$ onto $\A_m(\mathcal LO)$, so the Bell correlation of
$\A_m(O_1)$ and $\A_m(O_2+Le)$ equals that of $\A_m(\mathcal LO_1)$
and $\A_m(\mathcal LO_2+(0,L\eone))$. Choose $R$ with
$\mathcal LO_1\cup\mathcal LO_2\subset D_R$. By isotony,
$O\subset O'$ implies $\A_m(O)\subset\A_m(O')$, the two algebras
lie in $\A_m(D_R)$ and $\A_m(D_R+(0,L\eone))$, which commute for
$L>2R$, and enlarging both algebras can only increase the supremum
\eqref{eq:beta}; so the Bell correlation is at most $\beta_m(R,L)=1$
for $L\ge C_*R$, and it is always at least $1$.
\end{proof}

\appendix

\section{Proof of the probability criterion}\label{app:probability}

This appendix proves Theorem~\ref{thm:probability}.

\emph{Step 1: reduction to reflections.} The choice
$A_i=B_j=I$ gives the value $1$
in \eqref{eq:beta}, so $\beta(\omega;\A,\B)=1$ unless some self-adjoint
contractions $A_0,A_1\in\A$ and $B_0,B_1\in\B$ give the CHSH expression
an absolute value larger than $1$. The expression is real, because the
$A_i$ and $B_j$ commute, and negating both observables in one of the
two regions
changes its sign, so we may assume a value larger than $1$. Every
self-adjoint contraction is a norm-convergent average of reflections,
$A=\frac12\int_{-1}^1(2\mathbf 1_{\{A\ge u\}}-I)\dd u$, where
$\mathbf 1_{\{A\ge u\}}$ is the spectral projection of $A$ for the
interval $[u,1]$, and these reflections lie in every von Neumann algebra that
contains $A$. The CHSH expression is
affine and norm-continuous in each of its four arguments, so a value
larger than $1$ at $A_0$ gives a value larger than $1$ at one of these
reflections; replacing $A_0,A_1,B_0,B_1$ in turn, a strict violation by
contractions becomes a strict violation by four reflections. For
reflections, the \emph{score} is
$S=\frac12\omega\bigl(A_0(B_0+B_1)+A_1(B_0-B_1)\bigr)$.

\emph{Step 2: cells and their symmetries.} With
$E_i^\pm=(I\pm A_i)/2$ and $P_j^\pm=(I\pm B_j)/2$, which are
projections, the sixteen \emph{cells} are the joint probabilities
$\omega(E_i^\pm P_j^\pm)$; the index pair $(i,j)$ is called the
\emph{setting pair}. A cell of the setting pair $(i,j)$ is
\emph{winning} when the product of its two outcomes equals the sign
with which $\omega(A_iB_j)$ enters
$S=\frac12(\omega(A_0B_0)+\omega(A_0B_1)+\omega(A_1B_0)-\omega(A_1B_1))$;
so the winning cells are $E_i^+P_j^+$ and $E_i^-P_j^-$ for
$(i,j)\ne(1,1)$, and $E_1^+P_1^-$ and $E_1^-P_1^+$ for $(i,j)=(1,1)$.
Since $\pm\omega(A_iB_j)$, with the sign it carries in $S$, equals the
sum of the two winning cells of $(i,j)$ minus the sum of the two losing
ones, and the four cells of each setting pair add up to $1$, the value
$S+2$ is the sum of the eight winning cells.
The three replacements
$(A_0,A_1,B_0,B_1)\mapsto(A_0,-A_1,B_1,B_0)$,
$(A_1,A_0,B_0,-B_1)$ and $(-A_0,-A_1,-B_0,-B_1)$ carry reflection
quartets to reflection quartets, leave $S$ unchanged, and permute the
eight winning cells, and every winning cell can be moved to every
other by composing them: the first two exchange the settings
at one site and flip one outcome to compensate the sign of the
$(1,1)$ correlator, and the third exchanges the two winning cells of
each setting pair. We may therefore assume that a least likely winning
cell is $w=\omega(E_0^-P_0^-)$.

\emph{Step 3: the probability form of $S>1$.} Write
$e_A=\omega(E_0^-)$, $e_B=\omega(P_0^-)$, $x_A=\omega(E_1^-)$,
$x_B=\omega(P_1^-)$. Since $A_i=I-2E_i^-$ and $B_j=I-2P_j^-$, and
$\omega((I-2E)(I-2P))=1-2\omega(E)-2\omega(P)+4\omega(EP)$, the four
correlators are $1-2e_A-2e_B+4w$, $1-2e_A-2x_B+4\omega(E_0^-P_1^-)$,
$1-2x_A-2e_B+4\omega(E_1^-P_0^-)$ and $1-2x_A-2x_B+4\omega(E_1^-P_1^-)$.
Inserting them into $S$, the terms in $x_A,x_B$ cancel, and using
$\omega(E_0^-P_1^+)=e_A-\omega(E_0^-P_1^-)$ and
$\omega(E_1^+P_0^-)=e_B-\omega(E_1^-P_0^-)$,
\[
S=1+2\bigl\{w-\omega(E_0^-P_1^+)-\omega(E_1^+P_0^-)-\omega(E_1^-P_1^-)\bigr\}.
\]
So $S>1$ reads
\begin{equation}
w>\omega(E_0^-P_1^+)+\omega(E_1^+P_0^-)+\omega(E_1^-P_1^-).
\label{eq:probability-CHSH}
\end{equation}
The three terms on the right are nonnegative, so $w>0$ and each of them
is smaller than $w$.

\emph{Step 4: the neighbouring winning cells.} The cells $E_0^-P_1^-$
and $E_1^-P_0^-$ are winning, so their probabilities are at least
$w$; by Step~3 they are also $e_A-\omega(E_0^-P_1^+)>e_A-w$ and
$e_B-\omega(E_1^+P_0^-)>e_B-w$. Hence
$\omega(E_0^-P_1^-)\ge\max\{w,e_A-w\}\ge e_A/2$ and
$\omega(E_1^-P_0^-)\ge e_B/2$.

\emph{Step 5: the marginals $x_A,x_B$.} With the exponents $a,b$ of
\eqref{eq:probability-exponents}, abbreviate $\overline k=K_{r_+,+}^{1/r_+}$ and
$\underline k=K_{r_-,-}^{-1/(r_--1)}$, so that \eqref{eq:event-orders} reads
$\omega(EP)\le\overline k(\omega(E)\omega(P))^{b}$ and
$\omega(EP)\ge\underline k(\omega(E)\omega(P))^{1/a}$. The upper bound at
the two neighbouring cells gives
$e_A/2\le\overline k(e_Ax_B)^b$ and $e_B/2\le\overline k(x_Ae_B)^b$. Since
$e_A\ge w>0$ and $e_B\ge w>0$, dividing and taking the power $1/b$,
\[
x_B\ge(2\overline k)^{-1/b}e_A^{(1-b)/b},\qquad x_A\ge(2\overline k)^{-1/b}e_B^{(1-b)/b}.
\]

\emph{Step 6: a lower bound on $w$.} The upper bound at $E_0^-P_0^-$
reads $w\le\overline k(e_Ae_B)^b$, that is, $e_Ae_B\ge(w/\overline k)^{1/b}$. By
\eqref{eq:probability-CHSH}, the lower bound in
\eqref{eq:event-orders} at $E_1^-P_1^-$, Step~5, and this inequality
raised to the positive power $(1-b)/(ab)$,
\begin{align*}
w>\omega(E_1^-P_1^-)&\ge\underline k(x_Ax_B)^{1/a}\\
&\ge\underline k(2\overline k)^{-2/(ab)}(e_Ae_B)^{(1-b)/(ab)}\\
&\ge\underline k\,2^{-2/(ab)}\overline k^{\,-(1+b)/(ab^2)}w^{(1-b)/(ab^2)},
\end{align*}
the exponents of $\overline k$ adding up to $-2/(ab)-(1-b)/(ab^2)=-(1+b)/(ab^2)$.
Dividing by $w^{(1-b)/(ab^2)}$ and using
\[
\overline k^{\,-(1+b)/(ab^2)}=K_{r_+,+}^{-\nu_+}\quad\text{and}\quad
\underline k=K_{r_-,-}^{-1/(r_--1)}=K_{r_-,-}^{-\nu_-},
\]
which hold by
\eqref{eq:probability-exponents} because $r_--1=a/(1-a)$,
\begin{equation}
w^\eta K_{r_+,+}^{\nu_+}K_{r_-,-}^{\nu_-}>2^{-2/(ab)}.
\label{eq:necessary-cell}
\end{equation}

\emph{Step 7: an upper bound on $w$ from clustering.} Let
$S_{\mathrm{prod}}$ be $S$ with each $\omega(A_iB_j)$ replaced by
$\omega(A_i)\omega(B_j)$. Since $\abs{\omega(A_i)}\le1$ and
$\frac12(\abs{x+y}+\abs{x-y})=\max\{\abs x,\abs y\}$,
$S_{\mathrm{prod}}\le\max_j\abs{\omega(B_j)}$. The
bound \eqref{eq:additive-cluster} on the four pairs gives
$\abs{S-S_{\mathrm{prod}}}\le2q_{\mathrm{cl}}$, so
$1<S\le\max_j\abs{\omega(B_j)}+2q_{\mathrm{cl}}$. For the $j$
attaining the maximum, the less probable outcome of $B_j$ has
probability $(1-\abs{\omega(B_j)})/2<q_{\mathrm{cl}}$. If that outcome
is $P_j^\pm$,
the cell $E_0^\pm P_j^\pm$ is winning, because $\omega(A_0B_j)$ enters
$S$ with the sign $+$, and $E_0^\pm P_j^\pm\le P_j^\pm$; so
$w\le\omega(P_j^\pm)<q_{\mathrm{cl}}$.

\emph{Conclusion.} Since $\eta>0$, $0<w<q_{\mathrm{cl}}$ and
$K_{r_\pm,\pm}\ge1$, we get
$w^\eta K_{r_+,+}^{\nu_+}K_{r_-,-}^{\nu_-}
<q_{\mathrm{cl}}^\eta K_{r_+,+}^{\nu_+}K_{r_-,-}^{\nu_-}\le2^{-2/(ab)}$
by \eqref{eq:probability-cutoff}, contradicting
\eqref{eq:necessary-cell}. Hence no reflection quartet has $S>1$, and
$\beta(\omega;\A,\B)=1$ by Step~1. \qed

\section{Proofs of the Gaussian comparison and of its consequences}
\label{app:gaussian}

\emph{Gaussian conventions.---}For $n$ modes with commutation matrix
$J_n$, a centred Gaussian state $\rho_V$ with covariance $V$ has
characteristic function $\Tr(\rho_Ve^{\ii x^T(q,p)})=e^{-x^TVx/2}$,
$x\in\R^{2n}$, where $(q,p)=(\Phi(\sx_1),\dots,\Phi(\sx_{2n}))$ is the
column of canonical operators in a basis as in
Section~\ref{sec:field}, in the normalization where the oscillator
ground state has covariance $I/2$ \cite{SLW18}. For faithful $\sigma$ with covariance
$V_\sigma$, abbreviate $\mathsf T_\tau=V_\sigma^{[\tau]}$ and
$\mathsf L=(\mathsf T_\tau^2-I/4)^{1/2}$, and work in Williamson
normal-mode coordinates of $\sigma$, $V_\sigma=\bigoplus_j\upsilon_jI_2$:
the trace formula below is coordinate independent, but the covariance
formula holds only there. For
$\Theta=\sigma^{-\tau/2}\rho\,\sigma^{-\tau/2}$, the two Gaussian
sandwich formulas
\begin{equation}
\Tr\Theta=\frac{\Tr\sigma^{1-\tau}}{\sqrt{\det(I-\mathsf Z)}},\qquad
V_\Theta=-\mathsf T_\tau+\mathsf L(\mathsf T_\tau-V_\rho)^{-1}\mathsf L,
\label{eq:sandwich}
\end{equation}
hold whenever $V_\rho<\mathsf T_\tau$ (verified in Step~2 below),
where $V_\Theta$ is the covariance
of the Gaussian state $\Theta/\Tr\Theta$ and $\mathsf Z$ is the matrix
of Theorem~\ref{thm:gaussian}. The trace formula is that of
Seshadreesan, Lami and Wilde \cite[Theorem~22]{SLW18}, who compute
$\Qtil_r$ for Gaussian states from the opposite sandwich
$\rho^{1/2}\sigma^{-\tau}\rho^{1/2}$; it transfers because
$\rho^{1/2}\sigma^{-\tau}\rho^{1/2}=(\sigma^{-\tau/2}\rho^{1/2})^*(\sigma^{-\tau/2}\rho^{1/2})$
and $\Theta=(\sigma^{-\tau/2}\rho^{1/2})(\sigma^{-\tau/2}\rho^{1/2})^*$
have the same nonzero spectrum. The covariance formula is specific to
the sandwich used here; both follow from the coherent-state (Bargmann)
representation~\cite{Bargmann61}. By
\eqref{eq:mode-temperature}, the $j$th mode of $\sigma$ is
${(1-e^{-\beta_j})e^{-\beta_jN_j}}$, $N_j$ its number operator, and
$t^{N_j}$ acts on the Bargmann space of entire functions by scaling
the $j$th variable with $t$; so $\Theta$ has the integral kernel of
$\rho$ with both arguments scaled by
$\mathbb T=\bigoplus_je^{\tau\beta_j/2}I_2$, times
$\prod_j(1-e^{-\beta_j})^{-\tau}$. The kernel of $\rho$ is the exponential
of a quadratic form determined by $V_\rho$, and the scaling keeps it
Gaussian, so the trace and the characteristic function of $\Theta$ are
Gaussian integrals; they converge, and $\Theta$ is trace class,
exactly when $V_\rho<\mathsf T_\tau$, and a Schur complement yields
\eqref{eq:sandwich} through the mode identities
\[
\mathsf T_\tau-\tfrac12I=(\mathbb T^2-I)^{-1},\qquad
\mathsf L=\mathbb T\,(\mathbb T^2-I)^{-1}.
\]

\begin{proof}[Proof of Theorem~\ref{thm:gaussian}]
\emph{Step 1: coordinate independence.} For a real invertible $S$, a
positive $\Gamma$ and a real symmetric $\Xi$, the matrix
$(S\Gamma S^T)^{-1/2}S\Gamma^{1/2}$ is orthogonal, so
$(S\Gamma S^T)^{-1/2}\,S\Xi S^T\,(S\Gamma S^T)^{-1/2}$ is orthogonally
equivalent to $\Gamma^{-1/2}\Xi\Gamma^{-1/2}$. A real linear change of
the canonical coordinates replaces $V_\rho$, $V_\sigma$ and
$\Therm_\tau(V_\sigma)$ by $SV_\rho S^T$ and so on, so with
$\Gamma=\Therm_\tau(V_\sigma)$ and $\Xi=V_\rho-V_\sigma$ it replaces
$\mathsf Z$ by an orthogonally equivalent matrix. Hence $\norm{\mathsf Z}$,
$\norm{\mathsf Z}_1$ and the hypothesis are coordinate independent, as
is $\Qtil_r(\rho\|\sigma)$. We work in Williamson coordinates of
$\sigma$, in which $V_\sigma$, $\mathsf T_\tau$,
$\mathsf L$ and $\mathsf T_{1-\tau}=V_\sigma^{[1-\tau]}$ are all
diagonal.

\emph{Step 2: the condition $V_\rho<\mathsf T_\tau$.} With
$\delta=\norm{\mathsf Z}\le(4k_r)^{-1}\le\frac18$,
$\mathsf T_\tau-V_\rho=\Therm_\tau(V_\sigma)-(V_\rho-V_\sigma)
=\Therm_\tau(V_\sigma)^{1/2}(I-\mathsf Z)\Therm_\tau(V_\sigma)^{1/2}>0$, so \eqref{eq:sandwich}
applies.

\emph{Step 3: the perturbed covariance.} At $V_\rho=V_\sigma$ the
sandwich is $\sigma^{1-\tau}$, whose normalized state has covariance
$\mathsf T_{1-\tau}$, so \eqref{eq:sandwich} gives
$\mathsf T_{1-\tau}=-\mathsf T_\tau+\mathsf L\,\Therm_\tau(V_\sigma)^{-1}\mathsf L$.
Subtracting this from \eqref{eq:sandwich} and using Step~2,
\[
\begin{gathered}
V_\Theta-\mathsf T_{1-\tau}
=\mathsf P\bigl((I-\mathsf Z)^{-1}-I\bigr)\mathsf P,\\
\mathsf P=\mathsf L\,\Therm_\tau(V_\sigma)^{-1/2},\qquad
\mathsf P^2=\mathsf T_{1-\tau}+\mathsf T_\tau .
\end{gathered}
\]
Mode by mode, with $u=\artanh(1/(2\upsilon_j))$, the entries of
$\mathsf T_\tau$ and $\mathsf T_{1-\tau}$ are $\frac12\coth(\tau u)$
and $\frac12\coth((1-\tau)u)$ by \eqref{eq:thermal-power}; since
$x\coth x$ increases and $\coth$ decreases,
$\coth(\tau u)/\coth(({1-\tau})u)\le\max\{1,({1-\tau})/\tau\}$, so
$\mathsf P^2\le k_r\mathsf T_{1-\tau}$. The relative change of
the covariance of $\Theta$, the analogue of $\mathsf Z$ for $\Theta$,
is
\[
\begin{aligned}
\mathsf Z_\Theta&=\mathsf T_{1-\tau}^{-1/2}(V_\Theta-\mathsf T_{1-\tau})\mathsf T_{1-\tau}^{-1/2}\\
&=(\mathsf T_{1-\tau}^{-1/2}\mathsf P)\,\mathsf Z(I-\mathsf Z)^{-1}\,(\mathsf P\mathsf T_{1-\tau}^{-1/2}).
\end{aligned}
\]
Since the eigenvalues of $\mathsf Z$ lie in $[-\delta,\delta]$ and
$\norm{\mathsf T_{1-\tau}^{-1/2}\mathsf P}^2\le k_r$,
\begin{equation}
\norm{\mathsf Z_\Theta}\le\frac{k_r\delta}{1-\delta}\le\frac27,\qquad
\norm{\mathsf Z_\Theta}_1\le\frac{k_r}{1-\delta}\norm{\mathsf Z}_1 .
\label{eq:purity-perturbation}
\end{equation}

\emph{Step 4: the purity function.} For a \emph{physical} covariance
$V$, one with $V+\frac\ii2J_n\ge0$, with symplectic eigenvalues
$\upsilon_j$, let
$\mathscr P_r(V)=\log\Tr\rho_V^r=-\sum_j\log f_r(\upsilon_j)$ with
$f_r(x)=(x+\frac12)^r-(x-\frac12)^r$; the second expression holds mode
by mode, since a thermal mode is $(1-e^{-\beta})e^{-\beta N}$, $N$
the number operator, with $e^{-\beta}=(\upsilon-\frac12)/(\upsilon+\frac12)$
by \eqref{eq:mode-temperature}, and
$\Tr((1-e^{-\beta})e^{-\beta N})^r=1/f_r(\upsilon)$. One checks
$0\le xf_r'(x)/f_r(x)\le r$ for $x\ge\frac12$. The symplectic
eigenvalues are the positive eigenvalues of $\ii J_nV$; for a function
$F$ of the covariance write
$\dd F(V)[\dot V]=\frac{\dd}{\dd t}F(V+t\dot V)\big|_{t=0}$, $\dot V$
symmetric. First-order perturbation at
$\Upsilon=\bigoplus_j\upsilon_jI_2$ with distinct $\upsilon_j$ gives $\dd\upsilon_j[\dot V]=\frac12\Tr\dot V_{jj}$ for the
$j$th $2\times2$ diagonal block, whence
$\abs{\dd\mathscr P_r(\Upsilon)[\dot V]}
\le\frac r2\sum_j\abs{\Tr(\upsilon_j^{-1}\dot V_{jj})}
\le\frac r2\norm{\Upsilon^{-1/2}\dot V\Upsilon^{-1/2}}_1$. By the
invariance of $\mathscr P_r$ under $V\mapsto SVS^T$ with $S$
symplectic, and by Step~1, at every faithful
covariance
\begin{equation}
\abs{\dd\mathscr P_r(V)[\dot V]}\le\frac r2\norm{V^{-1/2}\dot VV^{-1/2}}_1 ,
\label{eq:purity-derivative}
\end{equation}
first where the $\upsilon_j$ are distinct, then everywhere by
continuity, since $\mathscr P_r(V)=-\frac12\Tr\log f_r(\abs{V^{1/2}\ii J_nV^{1/2}})$
is smooth on faithful covariances.

\emph{Step 5: integration.} The segment
$V(v)=\mathsf T_{1-\tau}^{1/2}(I+v\mathsf Z_\Theta)\mathsf T_{1-\tau}^{1/2}$,
$0\le v\le1$, joins $\mathsf T_{1-\tau}$ to $V_\Theta$ through physical
covariances, faithful for $v<1$ because
$\mathsf T_{1-\tau}+\frac\ii2J_n>0$. Its derivative is
$\mathsf T_{1-\tau}^{1/2}\mathsf Z_\Theta\mathsf T_{1-\tau}^{1/2}$, and by
Step~1 the relative perturbation $V(v)^{-1/2}\dot VV(v)^{-1/2}$ has the
trace norm of $\mathsf Z_\Theta(I+v\mathsf Z_\Theta)^{-1}$, at most
$\norm{\mathsf Z_\Theta}_1/(1-\norm{\mathsf Z_\Theta})$. Integrating
\eqref{eq:purity-derivative} and letting $v\to1$,
\[
\abs{\mathscr P_r(V_\Theta)-\mathscr P_r(\mathsf T_{1-\tau})}
\le\frac r2\frac{\norm{\mathsf Z_\Theta}_1}{1-\norm{\mathsf Z_\Theta}}
\le\frac r2\frac{k_r\norm{\mathsf Z}_1}{(1-\delta)(1-\norm{\mathsf Z_\Theta})}.
\]

\emph{Step 6: assembly.} Since $\Theta/\Tr\Theta$ is Gaussian with
covariance $V_\Theta$, \eqref{eq:renyi} and \eqref{eq:sandwich} give
$\Qtil_r(\rho\|\sigma)=(\Tr\Theta)^re^{\mathscr P_r(V_\Theta)}
=(\Tr\sigma^{1-\tau})^r\det(I-\mathsf Z)^{-r/2}e^{\mathscr P_r(V_\Theta)}$,
and $(\Tr\sigma^{1/r})^re^{\mathscr P_r(\mathsf T_{1-\tau})}=\Tr\sigma=1$
because $\sigma^{1/r}/\Tr\sigma^{1/r}$ has covariance
$\mathsf T_{1-\tau}$ and $r(1-\tau)=1$. Hence
\[
\log\Qtil_r(\rho\|\sigma)
=-\frac r2\Tr\log(I-\mathsf Z)+\mathscr P_r(V_\Theta)-\mathscr P_r(\mathsf T_{1-\tau}).
\]
The bound $\abs{\log(1-t)}\le\abs t/(1-\abs t)$ on the eigenvalues of
$\mathsf Z$ gives $\abs{\Tr\log(I-\mathsf Z)}\le\norm{\mathsf Z}_1/(1-\delta)$.
With Step~5, $1/(1-\delta)\le8/7$ and $1/(1-\norm{\mathsf Z_\Theta})\le7/5$,
\[
\log\Qtil_r(\rho\|\sigma)
\le r\Bigl(\frac47+\frac45k_r\Bigr)\norm{\mathsf Z}_1
\le2rk_r\norm{\mathsf Z}_1 ,
\]
the last step because $k_r\ge2$.
\end{proof}

\begin{proof}[Proof of Lemma~\ref{lem:renyi-consequences}]
Both parts are known: \eqref{eq:effect-renyi} is the two-outcome case
of the data-processing inequality for the sandwiched divergence, proved
for all $r>1$ by Frank and Lieb~\cite{FL13} and by Beigi~\cite{Beigi13}
and used in this form in strong-converse arguments \cite{WWY14}, and at
$r=2$ the quantity $\Qtil_2(\rho\|\sigma)-1$ is a quantum
$\chi^2$-divergence, which dominates the squared trace distance
\cite{TKRVV10}. We include the short proofs.

\emph{The effect bound.} Finiteness of $\Qtil_r(\rho\|\sigma)$ means
that $\Theta=\sigma^{-\tau/2}\rho\sigma^{-\tau/2}$ has finite $r$-norm
$\norm\Theta_r=\Qtil_r(\rho\|\sigma)^{1/r}$, and
$\rho=\sigma^{\tau/2}\Theta\sigma^{\tau/2}$. Since $\tau+1/r=1$,
H\"older's inequality for Schatten norms,
$\abs{\Tr(XY)}\le\norm X_p\norm Y_q$ for $1/p+1/q=1$ \cite{Simon05},
gives
$\rho(M)=\Tr(\Theta\,\sigma^{\tau/2}M\sigma^{\tau/2})
\le\norm\Theta_r\norm{\sigma^{\tau/2}M\sigma^{\tau/2}}_{1/\tau}$. The
operator $\sigma^{\tau/2}M\sigma^{\tau/2}$ has the same nonzero spectrum
as $M^{1/2}\sigma^\tau M^{1/2}$, and Hansen's inequality for the operator
monotone function $t\mapsto t^\tau$ (monotone: $0\le X\le Y$ implies
$X^\tau\le Y^\tau$) and the contraction $M^{1/2}$ \cite{Hansen80}
gives
$M^{1/2}\sigma^\tau M^{1/2}\le(M^{1/2}\sigma M^{1/2})^\tau$; so the
second factor is at most $(\Tr M^{1/2}\sigma M^{1/2})^\tau=\sigma(M)^\tau$.
Raising $\rho(M)\le\Qtil_r^{1/r}\sigma(M)^\tau$ to the power $r$ gives
\eqref{eq:effect-renyi}.

\emph{Comparison of the orders $2$ and $r$.} For $r\ge2$, interpolation by
Hadamard's three-lines theorem for Schatten norms \cite{Simon05},
applied to $z\mapsto\sigma^{-\tau z/2}\rho\sigma^{-\tau z/2}$ on
the strip $0\le\Rea z\le1$, between the trace norm at $\Rea z=0$
(where the sandwiching factors are unitary) and the $r$-norm at
$\Rea z=1$, evaluated at $z=1/(2\tau)\in(0,1]$ where the interpolated
exponent is $2$, gives
$\Qtil_2(\rho\|\sigma)\le\Qtil_r(\rho\|\sigma)^{1/(r-1)}$. To justify
the interpolation for the unbounded inverse powers, restrict
$\rho,\sigma$ to the spectral subspace of $\sigma$ for $[\epsilon,1]$,
where $\sigma^{-1}$ is bounded, and let $\epsilon\downarrow0$. This is the monotonicity of the sandwiched
divergence in its order~\cite[Theorem~7]{Beigi13},
\cite[Theorem~7]{ML13}.

\emph{The trace distance.} With
$\Theta_{1/2}=\sigma^{-1/4}\rho\sigma^{-1/4}$, the sandwich $\Theta$ at
$r=2$, i.e.\ $\tau=\frac12$,
one has $\Tr(\Theta_{1/2}\sigma^{1/2})=\Tr\rho=1$, so
$\norm{\Theta_{1/2}-\sigma^{1/2}}_2^2=\Qtil_2(\rho\|\sigma)-1$, and
$\rho-\sigma=\sigma^{1/4}({\Theta_{1/2}-\sigma^{1/2}})\sigma^{1/4}$ with
$\norm{\sigma^{1/4}}_4=1$. H\"older's inequality gives
$\norm{\rho-\sigma}_1\le\sqrt{\Qtil_2(\rho\|\sigma)-1}$, and the order
comparison gives \eqref{eq:trace-distance}.
\end{proof}

\section{Thermal weights and compression}\label{app:compression}

\begin{proof}[Proof of Lemma~\ref{lem:thermal-comparison}]
By the functional calculus of $d=\abs{\tfrac\ii2\mathsf J}$, as in
\eqref{eq:reserve-finite}, it suffices to show that
$\mathfrak q_\tau(d)=(d\coth(\tau\artanh d)-1)(1-d^2)^{-\tau}\ge c_\tau$
for $0<d<1$; at $d=0$ and $d=1$, which the spectrum may reach on the
full source space, \eqref{eq:thermal-d} means the continuous extension
(values $(1-\tau)/\tau$ and $0$), and the bound follows by continuity.
Write $d=\tanh u$, $u>0$ (on mode $j$ of a finite system,
$u=\beta_j/2$). Then $(1-d^2)^{-\tau}=\cosh(u)^{2\tau}$
and $\tanh u\coth(\tau u)-1=\sinh((1-\tau)u)/(\cosh u\sinh(\tau u))$,
so
\[
\mathfrak q_\tau(\tanh u)=\frac{\sinh((1-\tau)u)}{\sinh(\tau u)}\cosh(u)^{2\tau-1},
\]
with limits $(1-\tau)/\tau$ at $u\to0$ and $2^{1-2\tau}$ at
$u\to\infty$. For fixed $u$ the function $\mathfrak h(v)=v\coth(vu)$ is
convex in $v$, with $\mathfrak h'(v)\le\mathfrak h'(1)\le\tanh u$ for
$0<v\le1$, the last inequality being $\tanh u\le u$. By the mean value
theorem,
$\frac{\dd}{\dd u}\log\mathfrak q_\tau(\tanh u)
=\mathfrak h(1-\tau)-\mathfrak h(\tau)+(2\tau-1)\tanh u
=(1-2\tau)(\mathfrak h'(v_0)-\tanh u)$ at some
$v_0$ between $\tau$ and $1-\tau$, which is $\le0$ for $\tau\le\frac12$ and $\ge0$ for
$\tau\ge\frac12$. So $\mathfrak q_\tau(\tanh u)$ is monotone in $u$, and its infimum is
the limit at infinity for $\tau\le\frac12$ and the limit at zero for
$\tau\ge\frac12$, which is $c_\tau$.
\end{proof}

\begin{proof}[Proof of Lemma~\ref{lem:compression}]
\emph{Step 1: compression of parallel sums and geometric means.} For
positive invertible $A,B$, the \emph{parallel sum}
$(A^{-1}+B^{-1})^{-1}$ has the variational formula
$\ip z{(A^{-1}+B^{-1})^{-1}z}=\inf_{x+y=z}\{\ip xAx+\ip yBy\}$, which shows
that restricting $x,y$ to the range of $\iota$ increases the infimum,
so $\iota^*(A^{-1}+B^{-1})^{-1}\iota\le((\iota^*A\iota)^{-1}+(\iota^*B\iota)^{-1})^{-1}$,
and likewise for more than two summands. The weighted geometric mean
$A\#_\tau B=A^{1/2}(A^{-1/2}BA^{-1/2})^\tau A^{1/2}$ has the integral
representation
$\frac{\sin\pi\tau}\pi\int_0^\infty v^{\tau-1}(A^{-1}+vB^{-1})^{-1}\dd v$
through parallel sums, so
$\iota^*(A\#_\tau B)\iota\le(\iota^*A\iota)\#_\tau(\iota^*B\iota)$
\cite{KA80}.

\emph{Step 2: the weight operator.} $G$ and $\overline G$ commute, before and
after compression, because their sum is $2I$ and
$\iota^*(2I)\iota=2I_F$. For commuting operators
$A\#_{1/2}B=(AB)^{1/2}$, so Step~1 at $\tau=\frac12$ gives
$\Res_F\ge\iota^*\Res\iota$. Since $x\mapsto x^{2a}$ is operator monotone ($0\le X\le Y$ implies
$X^{2a}\le Y^{2a}$) and operator concave on $[0,\infty)$ for
$0<2a\le1$,
$\Res_F^{2a}\ge(\iota^*\Res\iota)^{2a}\ge\iota^*\Res^{2a}\iota$, the
second step being Jensen's inequality
$\iota^*X^{2a}\iota\le(\iota^*X\iota)^{2a}$ for the isometry $\iota$~\cite{Hansen80}.
This is
\eqref{eq:low-compression}.

\emph{Step 3: the thermal weight.} For an integer $r\ge2$, the thermal
increment is a parallel sum of geometric means,
\begin{equation}
\Therm_{(r-1)/r}(I)=\Bigl[\sum_{k=1}^{r-1}(G\#_{k/r}\overline G)^{-1}\Bigr]^{-1}
\label{eq:parallel-thermal}
\end{equation}
(for noninvertible $G,\overline G$ both sides are decreasing limits
over $G+\epsilon I$, $\overline G+\epsilon I$; Step~4). Indeed, at a
spectral value $d=\tanh u$ of $\frac\ii2\mathsf J$---here $d$ is
signed, but both sides of \eqref{eq:parallel-thermal} are even in it,
so $u>0$ suffices---one has
$G=1+d$, $\overline G=1-d$ and
$G\#_{k/r}\overline G=(1+d)e^{-2uk/r}$, so the
right-hand side equals
\[
\frac{(1+d)e^{-2u(r-1)/r}\bigl(1-e^{-2u/r}\bigr)}{1-e^{-2u(r-1)/r}}
=\frac{\sinh(u/r)}{\cosh u\,\sinh((r-1)u/r)},
\]
which is
$d\coth(\frac{r-1}r\artanh d)-1$ by the computation in the previous
proof. Compressing each geometric mean and then the parallel sum by
Step~1, and using \eqref{eq:parallel-thermal} on $F$, whose own $G$,
$\overline G$ are $G_F,\overline G_F$, gives
$\Therm_b(I_F)\ge\iota^*\Therm_b(I)\iota$, and
Lemma~\ref{lem:thermal-comparison} at $\tau=(r-1)/r\ge\frac12$, where
$c_\tau=(1-\tau)/\tau=(r-1)^{-1}$, gives the second inequality in
\eqref{eq:high-compression}.

\emph{Step 4: the full source space.} On the full space $G$ and
$\overline G$ are injective but need not be invertible. The
inequalities hold for $G+\epsilon I$ and $\overline G+\epsilon I$, and
as $\epsilon\downarrow0$ the geometric means and parallel sums decrease
to those of $G,\overline G$, which is how the noninvertible case is
defined, while the compressions converge to the invertible $G_F$,
$\overline G_F$. The inequalities pass to the limit.
\end{proof}

\section{Complex-shell evaluation and weighted moments}
\label{app:evaluation}

\begin{proof}[Proof of Lemma~\ref{lem:standard}]
The components of an element of $\Hh_{\mu,U}$ are distributions
supported in $\overline U$, with $H^{\mp1/2}$ Sobolev norms (Fourier
weights $(1+\abs\bp^2)^{\mp1/2}$) controlled by the covariance norm
for fixed $\mu$; their Fourier transforms therefore extend to entire
functions on $\C^3$, by the Paley--Wiener--Schwartz theorem
\cite{Hormander03}. If $\kappa_U(f,g)=0$, then
$\hat f(\bp)+\ii\sqrt{\mu^2+\abs\bp^2}\,\hat g(\bp)=0$ on $\R^3$; for
fixed real $\bp_\perp=(p_2,p_3)$, squaring and continuing in the first
variable gives
$\hat f(z_1,\bp_\perp)^2
=-(z_1^2+\mu^2+\abs{\bp_\perp}^2)\,\hat g(z_1,\bp_\perp)^2$
for all $z_1\in\C$. At $z_1=\ii\sqrt{\mu^2+\abs{\bp_\perp}^2}$ the
right side has a simple zero, so it has odd vanishing order unless it
vanishes identically, while the left side has even order; both sides
vanish, and $f=g=0$. For the density, in the spatial realization
$\kappa_U(f,g)=2^{-1/2}\bigl((\mu^2-\nabla^2)^{-1/4}f+\ii(\mu^2-\nabla^2)^{1/4}g\bigr)$ of the
one-particle space as $L^2(\R^3)$, a vector $v$ orthogonal to the
complex range gives $u=(\mu^2-\nabla^2)^{-1/4}v$ with $u=0$ and $(\mu^2-\nabla^2)^{1/2}u=0$ on
$U$. The function $W(t)=e^{-t(\mu^2-\nabla^2)^{1/2}}u$, with $t>0$
a Euclidean time, solves
the elliptic equation $(-\partial_t^2-\nabla^2+\mu^2)W=0$ with
Dirichlet and Neumann data $u$ and $-(\mu^2-\nabla^2)^{1/2}u$ at
$t=0$, which vanish on a ball
$B\subset U$; its extension by zero across $\{0\}\times B$ is a weak
solution on the connected set $((0,\infty)\times\R^3)\cup(\R\times B)$
that vanishes on an open subset, hence everywhere by analytic elliptic
regularity and unique continuation \cite{Hormander03}. So $u=0$ and
$v=0$. For regions with smooth boundary this is the anti-locality
argument of Figliolini and Guido \cite{FG89}.
\end{proof}

\emph{The modular operator of a ball.---}Let
$\mathcal K_U=\kappa_U\mathcal X_{\mu,U}$ be the closed real subspace of
one-particle vectors of real sources in $U$. By
Lemma~\ref{lem:standard}, $\mathcal K_U\cap\ii\mathcal K_U=\{0\}$ and
$\mathcal K_U+\ii\mathcal K_U$ is dense, so $\mathsf s_U(u+\ii v)=u-\ii v$
and $\Delta_U=\mathsf s_U^*\mathsf s_U$ are the standard-subspace
operators of Section~\ref{sec:estimates}, here for
$\mathcal K=\mathcal K_U$.
In the polar decomposition $\kappa_U=\mathcal U_UG_U^{1/2}$ the factor
$\mathcal U_U$ is unitary. The defining relation
$\mathsf s_U\kappa_U\sx=\kappa_U\mathsf c_U\sx$, where $\mathsf c_U$ is
complex conjugation on $\Hh_{\mu,U}$, and
$\overline G_U=\mathsf c_UG_U\mathsf c_U$ give
$\norm{\Delta_U^{1/2}\kappa_U\sx}^2=\norm{\kappa_U\mathsf c_U\sx}^2
=\ip\sx{\overline G_U\sx}$; comparing quadratic forms on the common
form domain $\Ran G_U^{1/2}$ (the set where both quadratic forms are
finite),
\begin{equation}
\Delta_U=\mathcal U_U\,\overline G_UG_U^{-1}\,\mathcal U_U^*,\qquad
\Delta_U^{\,c}\kappa_U=\mathcal U_U\,G_U^{1/2-c}\overline G_U^{\,c}
\label{eq:modular-source}
\end{equation}
for $0\le c\le\frac12$ \cite{RvD77,LS16}.

The proof of \eqref{eq:inclusion} given in
Section~\ref{sec:estimates} applies verbatim to any standard subspace
$\mathcal K\supset\mathcal K_U$, in particular to that of a Lorentz
transformed wedge; equivalently, $\Delta_{\mathcal K}^c\Delta_U^{-c}$
is a contraction \cite[Section~2]{BDL07}.

\emph{The wedge identity \eqref{eq:wedge-contour}.---}In rapidity
coordinates $p_0=m_\perp\cosh\chi$, $p_1=m_\perp\sinh\chi$, where
$\chi$ is the rapidity, $\bp_\perp=(p_2,p_3)$ and
$m_\perp=\sqrt{\mu^2+\abs{\bp_\perp}^2}$, the invariant measure is
$\dd\mathfrak m_\mu=(2(2\pi)^3)^{-1}\dd\chi\dd^2p_\perp$ and the boosts
along the first axis translate $\chi$. By the Bisognano--Wichmann
theorem \cite{BW75,BGL02}, $\Delta_{\mathcal W_0}^{\ii y}$ is the boost
of rapidity $2\pi y$, so after a Fourier transform in $\chi$ it is
multiplication by $e^{-2\pi\ii yk}$, $k$ the variable dual to $\chi$,
and $\Delta_{\mathcal W_0}^{\,c}$
is multiplication by $e^{-\varphi k}$, $\varphi=2\pi c$. By the Paley--Wiener
theorem for a strip \cite{Hormander03}, a wave function $\psi$ whose restriction to each
line of fixed $\bp_\perp$ extends holomorphically to $0<\Ima\chi<\varphi$,
with square-integrable bounds, lies in the domain of
$\Delta_{\mathcal W_0}^{\,c}$ and
$(\Delta_{\mathcal W_0}^{\,c}\psi)(\chi,\bp_\perp)=\psi(\chi+\ii\varphi,\bp_\perp)$.
For a smooth source $\sx$ with compact support in $\{x^1<0\}$, the
function $\chi\mapsto\Psi_\sx(\mathcal L_1(\chi)(m_\perp,0,\bp_\perp))$,
where $\mathcal L_1(\chi)$ is the boost of rapidity $\chi$ along the
first axis,
$(p_0,p_1)\mapsto(p_0\cosh\chi+p_1\sinh\chi,\,p_1\cosh\chi+p_0\sinh\chi)$,
is entire and equals $\kappa_U\sx$ at real $\chi$. On the strip, write
the continued rapidity as $\chi+\ii\varphi'$ with
$0\le\varphi'\le\varphi$; then $\abs{e^{-\ii\bz\cdot\bx}}\le1$ on the
support because $\Ima z_1=m_\perp\cosh\chi\sin\varphi'\ge0$, and integration by parts gives
rapid decay in $\abs\bz$. Hence
$\Delta_{\mathcal W_0}^{\,c}\kappa_U\sx=\Psi_\sx\circ b_\varphi$, since
$\mathcal L_1(\chi+\ii\varphi)=\mathcal L_1(\ii\varphi)\mathcal L_1(\chi)$. Translating
the wedge edge by $\ell$ conjugates $\Delta$ by the translation
$e^{-\ii\ell p_1}$, and the continuation $p_1\mapsto p_1\cos\varphi+\ii p_0\sin\varphi$
turns $e^{\ii\ell p_1}$ into a phase times $e^{-\ell p_0\sin\varphi}$,
which gives \eqref{eq:wedge-contour} for smooth sources in the base of
$\mathcal W_\ell$. For a proper orthochronous Lorentz transformation
$\mathcal L$, the subspace of $\mathcal L\mathcal W_\ell$ is
$\mathsf U(\mathcal L)\mathcal K_{\mathcal W_\ell}$, where $\mathsf U$
is the unitary one-particle representation of the Poincar\'e group,
$(\mathsf U(\mathcal L)\psi)(p)=\psi(\mathcal L^{-1}p)$; its modular
operator is the conjugate of $\Delta_{\mathcal W_\ell}$, and if
$D_T\subset\mathcal L\mathcal W_\ell$ then every real source in
$U\subset B_T(0)$ has $\kappa_U\sx\in\mathcal K_{\mathcal L\mathcal W_\ell}$%
---and so, by complex linearity, every $\sx\in\Hh_{\mu,U}$ has
$\kappa_U\sx\in\mathcal K_{\mathcal L\mathcal W_\ell}
+\ii\mathcal K_{\mathcal L\mathcal W_\ell}
=\operatorname{Dom}\mathsf s_{\mathcal L\mathcal W_\ell}$---%
with $\Psi_\sx(\mathcal Lb_\varphi(p))$ in place of $\Psi_\sx(b_\varphi(p))$
by Lorentz covariance of the field and the identity theorem on the
connected complex manifold $H_\mu^{\C}$. Finally, since
$\mathsf s_{\mathcal W_\ell}$ fixes $\mathcal K_{\mathcal W_\ell}$
pointwise,
$\norm{\Delta_{\mathcal W_\ell}^{1/2}\xi}=\norm\xi$ there, so
$\norm{\Delta^c_{\mathcal W_\ell}\xi}\le2\norm\xi$, and approximation
of an arbitrary $\sx\in\Hh_{\mu,U}$ by smooth sources, together with the
continuity of $\sx\mapsto\Psi_\sx(z)$ on $\Hh_{\mu,U}$ (the values
$\hat f(\bz)$, $\hat g(\bz)$ are pairings of $f\in H^{-1/2}$ and
$g\in H^{1/2}$ with a fixed smooth function), extends the identity to
all sources.

\emph{Averaged wedges.---}For $U\subset B_T(0)$ and a compact family
$\mathscr L$ of Lorentz transformations, the choice
$\ell_{\mathscr L}>\sqrt2\sup_{\mathcal L\in\mathscr L}\norm{\mathcal L^{-1}}_{\mathrm E}$,
with $\norm\cdot_{\mathrm E}$ the Euclidean operator norm on $\R^4$,
ensures $D_T\subset\mathcal L\mathcal W_{\ell_{\mathscr L}T}$ for every
$\mathcal L\in\mathscr L$, because $D_T$ lies in the Euclidean ball of
radius $T$ and $x^1+\abs{x^0}\le\sqrt2\abs x_{\mathrm E}$. The wedge
identity and \eqref{eq:inclusion}, integrated against a probability
measure $\vartheta$ on $\mathscr L$, give
\begin{equation}
\begin{split}
\int_{\mathscr L}\int_{H_\mu^+}&e^{-2\ell_{\mathscr L}Tp_0\sin\varphi}
\abs{\Psi_\sx(\mathcal Lb_\varphi(p))}^2\dd\mathfrak m_\mu(p)\dd\vartheta(\mathcal L)\\
&\le\norm{G_U^{1/2-c}\overline G_U^{\,c}\sx}^2 .
\end{split}
\label{eq:averaged-wedges}
\end{equation}

\begin{lemma}[Regular shell patches]\label{lem:regular-patches}
Let $0<c<\frac12$, $\varphi=2\pi c$, and let
$\mathscr Z=\{(t,z):0\le t\le1,\ z\in H_t^{\C},\ \abs z_{\mathrm E}\le16,\ (\Rea z)^2\ge16\}$,
where $(\Rea z)^2=\Rea z\cdot\Rea z$ is the Minkowski square, positive
for spacelike vectors. There are a compact set
$\mathscr L_c$ of proper orthochronous Lorentz transformations with
$\sup_{\mathscr L_c}\norm{\mathcal L^{-1}}_{\mathrm E}<32/\sin\varphi$, a
probability measure $\vartheta_c$ on $\mathscr L_c$ with smooth density
with respect to Haar measure, a
neighbourhood $\mathcal V$ of $\mathscr Z$ in the shell family
$\{(t,z):0\le t\le1,\ z\in H_t^{\C}\}$, and $\delta_c>0$, such that
the image of the measure
$\dd\vartheta_c\otimes\mathbf 1_{\{p_0\le M_c\}}\dd\mathfrak m_t$,
$M_c=32/\sin\varphi$, under $(\mathcal L,p)\mapsto\mathcal Lb_\varphi(p)$
assigns to every subset of the slice
$\mathcal V_t=\{z:(t,z)\in\mathcal V\}$ at least $\delta_c$ times its
six-dimensional Euclidean volume $\vol$; the formulas for
$\dd\mathfrak m_t$ extend to $t=0$, the measure on the forward light
cone. Moreover, there are a chart radius $r_c>0$ and
$C_{\mathrm{sub}}\ge1$ such that every $(t,z)\in\mathscr Z$ has a
holomorphic chart of radius $r_c$ in $H_t^{\C}$ centred at $z$, with
image $\mathcal B_z\subset\mathcal V_t$ and volume Jacobian between $1$
and $C_{\mathrm{sub}}$; hence every nonnegative plurisubharmonic $F$
on $\C^4$ (subharmonic on every complex line), in particular
the squared modulus of any holomorphic function, satisfies
$F(z)\le C_{\mathrm{sub}}\vol(\mathcal B_z)^{-1}\int_{\mathcal B_z}F\dd\vol$.
\end{lemma}

\begin{proof}[Proof (condensed)]
A point $z=u+\ii v\in H_t^{\C}$ with $(t,z)\in\mathscr Z$ has
$u\cdot v=0$, $u^2=\rho_0^2$ with $\rho_0\in[4,16]$, and
$v^2=t^2+\rho_0^2$, since $z\cdot z=u^2-v^2+2\ii\,u\cdot v=-t^2$. The
pair $(u,v)$ is the Lorentz image of the
real and imaginary parts of $b_\varphi(p)$ at $p_1=p_3=0$,
$p_0=\sqrt{t^2+\rho_0^2}/\sin\varphi$, $p_2=\sqrt{p_0^2-t^2}$, where
these parts are orthogonal with the same Minkowski squares; the matching transformation
$\mathcal L=\mathcal L_{\mathrm{tar}}\mathcal L_{\mathrm{src}}^{-1}$,
where $\mathcal L_{\mathrm{src}}$ and $\mathcal L_{\mathrm{tar}}$ are
proper orthochronous transformations carrying the standard frame to
orthonormal frames whose first two spacelike legs are the two pairs,
satisfies
$\norm{\mathcal L}_{\mathrm E}=\norm{\mathcal L^{-1}}_{\mathrm E}\le2\sqrt{34}/\sin\varphi<32/\sin\varphi$
(the Euclidean norms of a Lorentz matrix and of its inverse coincide) and
$p_0\le\sqrt{257}/\sin\varphi<32/\sin\varphi$: for the
normalized legs $u/\rho_0$ and $v/\sqrt{t^2+\rho_0^2}$, the bound
$\abs z_{\mathrm E}^2=2(u_0^2+v_0^2)+2\rho_0^2+t^2\le256$ gives
\[
\frac{u_0^2}{\rho_0^2}+\frac{v_0^2}{t^2+\rho_0^2}
\le\frac{u_0^2+v_0^2}{\rho_0^2}\le\frac{128}{\rho_0^2}-1\le7,
\]
so, completing the pair by a third spacelike leg with zero time
component, the Lorentz factor of $\mathcal L_{\mathrm{tar}}$ (the time
component of the image of the time axis) is at most $\sqrt8$; that of
$\mathcal L_{\mathrm{src}}$ is
$\sqrt{1+(17/16)\cot^2\varphi}\le\sqrt{17}/(4\sin\varphi)$, since
$t\le1$ and $\rho_0^2\ge16$; a proper orthochronous matrix of Lorentz
factor $\gamma$ has Euclidean norm
$\gamma+\sqrt{\gamma^2-1}\le2\gamma$, and the two norms multiply. The
map $(\mathcal L,p)\mapsto\mathcal Lb_\varphi(p)$ is a submersion at such
preimages: the Lorentz orbit of the ordered pair $(u,v)$ has dimension
five (the stabilizer of an orthogonal spacelike pair is the
one-parameter boost group of the complementary timelike plane), and
varying $p_2$ changes the invariant $(\Rea z)^2$ with
nonzero derivative:
\[
\begin{gathered}
(\Rea b_\varphi)^2=-p_0^2\cos^2\varphi+p_2^2
=(t^2+p_2^2)\sin^2\varphi-t^2,\\
\partial_{p_2}(\Rea b_\varphi)^2=2p_2\sin^2\varphi\ne0 .
\end{gathered}
\]
Including $t$ as a parameter, the submersion
theorem and Fubini's theorem give, near each point of $\mathscr Z$, a
push-forward density that is continuous in $(t,z)$, $t=0$ included,
and positive (the preimages have $p_2\ge\rho_0\ge4$, uniformly
away from the apex of the cone, so $\dd\mathfrak m_t$ and the local
Jacobians stay regular down to $t=0$); compactness of $\mathscr Z$ gives a finite cover, one
smooth Lorentz density supported in the open set
$\{\norm{\mathcal L^{-1}}_{\mathrm E}<32/\sin\varphi\}$, and $\delta_c>0$.
The charts come from the holomorphic implicit function theorem applied
to $z\cdot z=-t^2$ in a coordinate $z_k$ with $\abs{z_k}\ge4/\sqrt3$,
uniformly over the compact set $\mathscr Z$; the Jacobian bounds and
the sub-mean-value property of $F$ composed with the chart give the
last claim.
\end{proof}

\begin{lemma}[Null analytic discs]\label{lem:null-discs}
If $z\in H_\mu^{\C}$, $\Energy\ge\mu$ and $\abs z_{\mathrm E}\le2\Energy$,
there is an analytic disc $\zeta\mapsto z(\zeta)$ in $H_\mu^{\C}$,
$\abs\zeta\le1$, with $z(0)=z$, whose boundary points satisfy
$\abs{\Rea\bz(\zeta)}\ge6\Energy$, $\abs{\Rea z_0(\zeta)}\le2\Energy$,
$(\Rea z(\zeta))^2\ge32\Energy^2$ and $\abs{z(\zeta)}_{\mathrm E}<14\Energy$;
in particular $(\mu/\Energy,z(\zeta)/\Energy)\in\mathscr Z$.
\end{lemma}

\begin{proof}
Choose $\mathbf n\in\C^3$, $\mathbf n\ne0$, with $\bz\cdot\mathbf n=0$
and $\mathbf n\cdot\mathbf n=0$ (an \emph{isotropic} vector): the
first equation defines a complex subspace of dimension at least two,
and a quadratic form on such a subspace has a nonzero isotropic
vector. Normalizing, $\abs{\Rea\mathbf n}=\abs{\Ima\mathbf n}=1$ and
$\Rea\mathbf n\perp\Ima\mathbf n$, so $\abs{\Rea(\zeta\mathbf n)}=1$
for $\abs\zeta=1$. The disc $z(\zeta)=z+8\Energy\zeta(0,\mathbf n)$ lies
in $H_\mu^{\C}$ because $\bz\cdot\mathbf n=\mathbf n\cdot\mathbf n=0$,
and on its boundary $z_0$ is unchanged while
$\abs{\Rea\bz(\zeta)}\ge8\Energy-\abs{\Rea\bz}\ge6\Energy$; the
Minkowski-square bound follows, and
$\abs{z(\zeta)}_{\mathrm E}\le(2+8\sqrt2)\Energy<14\Energy$.
\end{proof}

\begin{proof}[Proof of Theorem~\ref{thm:evaluation}]
Let $\sx$ be smooth with compact support in $U$, and write
$\mu=t\Energy$ with $t\in(0,1]$; the substitution $p=\Energy p'$ scales
the shell measure by $\Energy^2$. For $(t,z)\in\mathscr Z$, the
plurisubharmonic function $F=\abs{\Psi_\sx(\Energy\,\cdot)}^2$, the chart
bound of Lemma~\ref{lem:regular-patches}, the push-forward bound on
$\mathcal B_z\subset\mathcal V_t$, the insertion of the factor
$e^{2\ell_cT(\Energy M_c-p_0)\sin\varphi}\ge1$ on $\{p_0\le\Energy M_c\}$, and
\eqref{eq:averaged-wedges} with $\mathscr L=\mathscr L_c$,
$\ell_{\mathscr L}=\ell_c=64/\sin\varphi$ (which exceeds
$\sqrt2\cdot32/\sin\varphi$, as required) give
\[
\abs{\Psi_\sx(\Energy z)}^2
\le\frac{C_{\mathrm{sub}}\Energy^{-2}e^{2\ell_cM_cT\Energy\sin\varphi}}
{\delta_c\vol(\mathcal B_z)}
\norm{G_U^{1/2-c}\overline G_U^{\,c}\sx}^2,
\]
with $\ell_cM_c\sin\varphi=2048/\sin(2\pi c)=\lambda_c$ and
$C_c=\max\{1,C_{\mathrm{sub}}/(\delta_c\inf_{\mathscr Z}\vol(\mathcal B_z))\}$,
finite because $\vol(\mathcal B_z)\ge\vol(B_{r_c})$. For a general $z$ with
$\abs z_{\mathrm E}\le2\Energy$, Lemma~\ref{lem:null-discs} and the mean
value property of holomorphic functions,
$\abs{\Psi_\sx(z)}\le\max_{\abs\zeta=1}\abs{\Psi_\sx(z(\zeta))}$,
reduce the claim to boundary points of the form $\Energy z'$ with
$(t,z')\in\mathscr Z$.
Density of smooth sources and continuity of both sides in $\sx$ give
the statement on $\Hh_{\mu,U}$. The set $U$ enters only through
$U\subset B_T(0)$.
\end{proof}

\emph{Weighted moments.---}For a multi-index $\alpha\in\N^3$ and a
centre $\bx_0$, let $\mathsf m_\alpha^f(\sx)=\int(\bx-\bx_0)^\alpha f(\bx)\dd^3x$
and $\mathsf m_\alpha^g(\sx)=\int({\bx-\bx_0})^\alpha g(\bx)\dd^3x$. For
$\frac12\le b<1$ fix
\begin{equation}
\Lambda_b=\frac{2^{15}}{\sin(\pi b)}=16\lambda_{b/2},\qquad
B_b=5\Lambda_b .
\label{eq:Lambda}
\end{equation}

\begin{lemma}[Weighted moments]\label{lem:moments}
Let $\frac12\le b<1$. There is a finite $M_b$, depending only on $b$,
such that for all $\mu,T>0$, all $\bx_0\in\R^3$, every nonempty bounded
open $U\subset B_T(\bx_0)$, every $\sx\in\Hh_{\mu,U}$ and every
$\alpha\in\N^3$,
\begin{equation}
\begin{aligned}
\abs{\mathsf m_\alpha^f(\sx)}
&\le M_b\frac{T}{1+\mu T}e^{\Lambda_b\mu T}(B_bT)^{\abs\alpha}
\norm{\Res_U^{\,b}\sx},\\
\abs{\mathsf m_\alpha^g(\sx)}
&\le M_b\frac{T^2}{(1+\mu T)^2}e^{\Lambda_b\mu T}(B_bT)^{\abs\alpha}
\norm{\Res_U^{\,b}\sx}.
\end{aligned}
\label{eq:moments}
\end{equation}
\end{lemma}

By the Riesz representation theorem, each bound says that there is a
vector $\mathsf u\in\Hh_{\mu,U}$, of norm at most the stated constant,
with $\mathsf m_\alpha^\bullet(\sx)=\ip{\mathsf u}{\Res_U^b\sx}$ for all
$\sx$.

\begin{proof}[Proof of Lemma~\ref{lem:moments}]
Work in coordinates centred at $\bx_0$; translation preserves
$\varsigma$, $\cov_\mu$ and $\Res_U$. By \eqref{eq:symmetric-evaluation},
for $\Energy\ge\mu$,
$\abs{\Psi_\sx(z)}\le\bigl(C_{(1-b)/2}^{1/2}+C_{b/2}^{1/2}\bigr)\Energy^{-1}e^{\lambda_{b/2}T\Energy}\norm{\Res_U^b\sx}$
on $\{z\in H_\mu^{\C}:\abs z_{\mathrm E}\le2\Energy\}$. For $\bz\in\C^3$ take
$\Energy=\mu+1/(\Lambda_bT)+\abs\bz$. The two shell points
$z_\pm=(\pm\sqrt{\mu^2+\bz\cdot\bz},\bz)$ have
$\abs{z_\pm}_{\mathrm E}<\sqrt2\Energy$, and
$\hat f(\bz)=\frac12(\Psi_\sx(z_+)+\Psi_\sx(z_-))$, so
$\abs{\hat f(\bz)}\le\bigl({C_{(1-b)/2}^{1/2}+C_{b/2}^{1/2}}\bigr)\Energy^{-1}e^{\lambda_{b/2}T\Energy}\norm{\Res_U^b\sx}$.
For $\hat g$, the difference ${\Psi_\sx(z_+)-\Psi_\sx(z_-)}=2\ii\sqrt{\mu^2+\bz\cdot\bz}\,\hat g(\bz)$
is useless near zero energy, so move $z_1$ on a circle: for
$\abs\zeta=1$ put $\bz_\zeta=(z_1+8\Energy\zeta,z_2,z_3)$. Then
$\abs{\mu^2+\bz_\zeta\cdot\bz_\zeta}\ge47\Energy^2$, the shell points
over $\bz_\zeta$ have Euclidean norm below $14\Energy$, and
\eqref{eq:symmetric-evaluation} at the energy $7\Energy$ gives
$\abs{\hat g(\bz_\zeta)}\le\bigl(C_{(1-b)/2}^{1/2}+C_{b/2}^{1/2}\bigr)(7\sqrt{47})^{-1}\Energy^{-2}e^{7\lambda_{b/2}T\Energy}\norm{\Res_U^b\sx}$;
the mean value property of $\zeta\mapsto\hat g(\bz_\zeta)$ transfers
the bound to $\hat g(\bz)$.
Since $\Lambda_b=16\lambda_{b/2}$, both exponents are at most
$\Lambda_bT\Energy=\Lambda_b\mu T+1+\Lambda_bT\abs\bz$, and
$\Energy^{-1}\le\Lambda_bT/(1+\mu T)$; so, for all $\bz\in\C^3$,
\[
\begin{gathered}
\abs{\hat f(\bz)}\le e\bigl(C_{(1-b)/2}^{1/2}+C_{b/2}^{1/2}\bigr)\Lambda_b\frac T{1+\mu T}\,
e^{\Lambda_bT(\mu+\abs\bz)}\norm{\Res_U^b\sx},\\
\abs{\hat g(\bz)}\le e\bigl(C_{(1-b)/2}^{1/2}+C_{b/2}^{1/2}\bigr)\Lambda_b^2\frac{T^2}{(1+\mu T)^2}\,
e^{\Lambda_bT(\mu+\abs\bz)}\norm{\Res_U^b\sx}.
\end{gathered}
\]
Finally $\partial^\alpha\hat f(0)=(-\ii)^{\abs\alpha}\mathsf m_\alpha^f(\sx)$,
and for $\abs\alpha\ge1$ the Cauchy estimate on the polydisc of radius
$\abs\alpha/(\sqrt3\Lambda_bT)$, where $e^{\Lambda_bT\abs\bz}\le e^{\abs\alpha}$,
gives the factor $\alpha!(\sqrt3\Lambda_bT/\abs\alpha)^{\abs\alpha}e^{\abs\alpha}
\le(e\sqrt3\Lambda_bT)^{\abs\alpha}\le(B_bT)^{\abs\alpha}$; for
$\alpha=0$ the displayed bounds at $\bz=0$ are \eqref{eq:moments}
directly. Thus
\eqref{eq:moments} holds with
$M_b=e\bigl(C_{(1-b)/2}^{1/2}+C_{b/2}^{1/2}\bigr)\Lambda_b^2$.
\end{proof}

\section{The forward and reverse estimates}\label{app:estimates}

Throughout this appendix, $C$ is a universal constant whose value may
change from occurrence to occurrence; subscripted constants keep their
values.

\emph{Covariance kernels.---}Away from the diagonal, the covariance
\eqref{eq:covariance} is given by integral kernels: for sources with
disjoint supports,
$\cov_\mu(\sx,\sy)=\iint[f(\bx)\mathfrak X_\mu(\by-\bx)f'(\by)
+g(\bx)\mathfrak P_\mu({\by-\bx})g'(\by)]\dd^3x\dd^3y$ with
\begin{equation}
\mathfrak X_\mu(\mathbf w)=\frac{\mu}{4\pi^2\varrho}K_1(\mu\varrho),\qquad
\mathfrak P_\mu(\mathbf w)=-\frac{\mu^2}{4\pi^2\varrho^2}K_2(\mu\varrho),
\label{eq:kernels}
\end{equation}
where $\varrho=\sqrt{\mathbf w\cdot\mathbf w}$ and $K_1,K_2$ are
modified Bessel functions of the second kind \cite{DLMF}. These formulas come from the
heat-kernel representation of $(\mu^2-\nabla^2)^{-1/2}$ and the integral
representation $K_1(x)=(x/4)\int_0^\infty t^{-2}e^{-t-x^2/(4t)}\dd t$
\cite[Eq.~10.32.10]{DLMF}, and $\mathfrak P_\mu=(-\nabla^2+\mu^2)\mathfrak X_\mu$
by the Bessel recurrences. From $K_j(x)=\int_0^\infty e^{-x\cosh u}\cosh(ju)\dd u$
\cite[Eq.~10.32.9]{DLMF}, for $x>0$,
\begin{equation}
xK_1(x)\le3(1+x)^{1/2}e^{-x},\qquad x^2K_2(x)\le8(1+x)^{3/2}e^{-x},
\label{eq:Bessel-scaled}
\end{equation}
and $\abs{K_j(z)}\le K_j(\Rea z)$ for $\Rea z>0$.

\begin{proof}[Proof of Theorem~\ref{thm:forward}]
Measure $\bx$ and $\by$ from the centres of the two balls, so that
$\mathbf w=s\eone+\by-\bx$, and put $d_{0,b}=4\sqrt3\,B_b$ and
$S_b=2(2\Lambda_b+2d_{0,b})=(4+80\sqrt3)\,\Lambda_b$, which is at
least $4$ and $4d_{0,b}$; let $s\ge S_b$ and $\xi=\mu s$.

\emph{Step 1: the kernels on a polydisc.} On the six-variable polydisc
(a product of discs, one per complex coordinate)
$\abs{x_j},\abs{y_j}\le2B_b$ one has
$\mathbf w\cdot\mathbf w=s^2(1+\epsilon)$ with
$\abs\epsilon\le2d_{0,b}/s+(d_{0,b}/s)^2\le9/16$ since
$s\ge4d_{0,b}$, so the principal branch
$\varrho=s\sqrt{1+\epsilon}$ is holomorphic there with
$\Rea\varrho\ge s-2d_{0,b}\ge s/2$ and $\abs\varrho\le5s/4$. With
\eqref{eq:Bessel-scaled}, the polydisc suprema of the two kernels are
at most $Cs^{-2}(1+\xi)^{1/2}e^{-\xi+2\mu d_{0,b}}$ and
$Cs^{-4}({1+\xi})^{3/2}e^{-\xi+2\mu d_{0,b}}$, and their Taylor coefficients
$c^f_{\alpha\beta},c^g_{\alpha\beta}$ in $\bx^\alpha\by^\beta$ are
bounded by these suprema times $(2B_b)^{-\abs\alpha-\abs\beta}$.
The series converge absolutely on the unit balls.

\emph{Step 2: the rank-one series.} For smooth real sources $\sx,\sy$
in the two balls, term-by-term integration gives
\[
\cov_\mu(\sx,\sy)=\sum_{\alpha,\beta}\bigl(c^f_{\alpha\beta}\mathsf m^f_\alpha(\sx)\mathsf m^f_\beta(\sy)
+c^g_{\alpha\beta}\mathsf m^g_\alpha(\sx)\mathsf m^g_\beta(\sy)\bigr).
\]
By Lemma~\ref{lem:moments} at $T=1$ and the Riesz representation,
there are, for $\bullet\in\{f,g\}$, vectors
$\mathsf u^\bullet_{A,\alpha}$, $\mathsf u^\bullet_{B,\beta}$ with
$\mathsf m^\bullet_\alpha(\sx)=\ip{\mathsf u^\bullet_{A,\alpha}}{\Res_A^b\sx}$,
$\mathsf m^\bullet_\beta(\sy)=\ip{\mathsf u^\bullet_{B,\beta}}{\Res_B^b\sy}$,
of norms at most $M_be^{\Lambda_b\mu}B_b^{\abs\alpha}/({1+\mu})$ for
$\bullet=f$ and $M_be^{\Lambda_b\mu}B_b^{\abs\alpha}/(1+\mu)^2$ for
$\bullet=g$. Put
$X_{b,\mu,s}=\sum_{\bullet,\alpha,\beta}c^\bullet_{\alpha\beta}\,
\mathsf u^\bullet_{A,\alpha}\ip{\mathsf u^\bullet_{B,\beta}}{\cdot}$.
Since $B_b/(2B_b)=\tfrac12$ and
$\sum_{\alpha,\beta}2^{-\abs\alpha-\abs\beta}=64$, the series converges
in trace norm, with
\[
\begin{split}
\norm{X_{b,\mu,s}}_1
\le{}& CM_b^2e^{(2\Lambda_b+2d_{0,b})\mu-\xi}\\
&\times\Bigl[\frac{s^{-2}(1+\xi)^{1/2}}{(1+\mu)^2}
+\frac{s^{-4}(1+\xi)^{3/2}}{(1+\mu)^4}\Bigr].
\end{split}
\]

\emph{Step 3: the constant and the identification.} Since
$s\ge S_b=2(2\Lambda_b+2d_{0,b})$, one has
$(2\Lambda_b+2d_{0,b})\mu\le\xi/2$, so the right side is at most
$2CM_b^2s^{-2}(1+\xi)^{3/2}e^{-\xi/2}\le F_bs^{-2}$ with
$F_b=2\cdot3^{3/2}e^{-1}CM_b^2$, the supremum of $(1+\xi)^{3/2}e^{-\xi/2}$
being attained at $\xi=2$. On pairs of smooth real sources,
$\ip\sx{\Res_A^bX_{b,\mu,s}\Res_B^b\sy}$ equals the series of Step~2,
hence $\cov_\mu(\sx,\sy)=\ip\sx{\mathsf C_{\mu,s}\sy}$; by sesquilinearity and
density the two bounded operators coincide.
\end{proof}

\emph{Smooth kernels in the joint norm.---}For $s\ge5$, let $K$ be a
kernel on $(B_2(0)\cup B_2(s\eone))^2$ with derivatives of order at
most two in each variable bounded by $M_0$. Fix a smooth function equal
to $1$ on $U_s$, supported within distance $1$ of it and with
$s$-independent derivative bounds, and let $K^\circ$ be $K$ multiplied
by that function in each variable. The Hilbert--Schmidt norm (the
square root of the sum of the squared singular values) of
$(1-\nabla^2)^{1/4}K^\circ(1-\nabla^2)^{1/4}$ is bounded by $CM_0$,
and the $L^2(\R^3)$ comparisons
$\norm{(1-\nabla^2)^{-1/4}f}_{L^2}
\le\max\{1,\mu\}^{1/2}\norm{({\mu^2-\nabla^2})^{-1/4}f}_{L^2}$ and,
for $g$ supported in $U_s$,
$\norm g_{L^2}^2\le2\norm{(\mu^2-\nabla^2)^{1/4}g}_{L^2}^2$ (less
than half of $\norm g_{L^2}^2$ sits at frequencies $\abs\bp<1$, by the
volume of $U_s$) give, for sources $\sx=(f,g)$, $\sy=(f',g')$ supported
in $U_s$ and in the joint norm,
\begin{equation}
\begin{gathered}
\Bigl|\iint fKf'\Bigr|\le C\max\{1,\mu\}M_0\norm\sx\norm\sy,\\
\Bigl|\iint gKg'\Bigr|\le CM_0\norm\sx\norm\sy ,
\end{gathered}
\label{eq:smooth-kernel}
\end{equation}
with $C$ independent of $s\ge5$ and $\mu$.

The polynomial \emph{selectors}
\begin{equation}
\mathsf p_N(z)=(1-z)^N\sum_{k=0}^{N-1}\binom{N+k-1}kz^k
\label{eq:selectors}
\end{equation}
have degree $2N-1$; $\mathsf p_N-1$ vanishes to order $N$ at $z=0$ and
$\mathsf p_N$ vanishes to order $N$ at $z=1$. So
$\mathsf p_N(x_1/s)\approx1$ near the ball $B_1(0)$ and
$\mathsf p_N(x_1/s)\approx0$ near $B_1(s\eone)$: multiplying a kernel
by such factors switches it on between the two balls and off within
each ball, as the kernel of $\JointD_{\mu,s}$ requires.

\begin{lemma}[Polynomial increments]\label{lem:polynomial}
There are universal finite constants
$C_1,C_2,c_{\mathrm{pol}},K_{\mathrm{pol}}$ and $S_{\mathrm{pol}}$ such
that for every $\mu>0$ and $s\ge S_{\mathrm{pol}}$, with $\xi=\mu s$,
there are real self-adjoint finite-rank operators $\mathsf F_n$ on
$\Hh_{\mu,U_s}$ with polynomial kernels, $\mathsf F_0=0$ and
$\mathsf F_n\to\JointD_{\mu,s}$ in operator norm, whose increments
$\mathsf G_n=\mathsf F_{n+1}-\mathsf F_n$ satisfy
\begin{align}
\norm{\mathsf G_n}_1&\le C_1(n+1)^3s^{-2}(1+\xi)^2e^{-\xi/2}(c_{\mathrm{pol}}/s)^n,
\label{eq:increment-unweighted}\\
\mathsf G_n&=\JointRes^{1/2}\mathsf B_n\JointRes^{1/2},\quad
\norm{\mathsf B_n}_1\le C_2e^{-\xi/2}e^{2\Lambda_0\mu T}K_{\mathrm{pol}}^n,
\label{eq:increment-weighted}
\end{align}
where $\Lambda_0=\Lambda_{1/2}=2^{15}$ and, here and below,
$B_{1/2},M_{1/2}$ are the constants of Lemma~\ref{lem:moments} at
$b=\frac12$, $T=s/2+1$ is the radius of a ball about $s\eone/2$
containing $U_s$ (the value at which Lemma~\ref{lem:moments} is
applied), and the $\mathsf B_n$ are real self-adjoint of finite rank.
\end{lemma}

\begin{proof}[Proof (condensed)]
By \eqref{eq:D-joint} and \eqref{eq:kernels}, the kernel of
$\JointD_{\mu,s}$ is $-\mathfrak X_\mu(\by-\bx)$ on field sources and
$-\mathfrak P_\mu(\by-\bx)$ on momentum sources when $\bx$ and $\by$
lie in different balls, and $0$ when they lie in the same ball. The approximating kernel $q_{n,\bullet}$ below is
a Taylor polynomial of the cross kernel, multiplied by selectors that
keep it only between the balls.
In the coordinates $\mathbf u=\bx/s$, $\mathbf v=\by/s$, with
components $u_j,v_j$, the cross kernels
$\mathfrak X_\mu(s(\mathbf v-\mathbf u))$ and
$\mathfrak P_\mu(s(\mathbf v-\mathbf u))$ are holomorphic
on the polydisc of radius $2r_0=1/50$ (so $r_0=1/100$) about
$(0,\eone)$, with suprema
$M_f=Cs^{-2}(1+\xi)^{1/2}e^{-\xi/2}$ and
$M_g=Cs^{-4}(1+\xi)^{3/2}e^{-\xi/2}$ by \eqref{eq:Bessel-scaled}, and
Taylor coefficients bounded by $M_\bullet r_0^{-\abs\alpha-\abs\beta}$.
Let $\mathcal T_{n,\bullet}$ be their Taylor polynomials of degree
$n-1$ in the six variables $(\mathbf u,\mathbf v-\eone)$ about the
origin, and define the polynomial kernel
\begin{align*}
q_{n,\bullet}(\bx,\by)={}&-\mathsf p_{2n}(u_1)\mathsf p_{2n}(1-v_1)\mathcal T_{n,\bullet}(\mathbf u,\mathbf v-\eone)\\
&-\mathsf p_{2n}(v_1)\mathsf p_{2n}(1-u_1)\mathcal T_{n,\bullet}(\mathbf v,\mathbf u-\eone),
\end{align*}
real and symmetric, of degree at most $9n$, with the selectors
\eqref{eq:selectors}; $\mathsf F_n$ is the operator with the
field-source form of $q_{n,f}$ plus the momentum-source form of
$q_{n,g}$. Its form depends on $\sx$ only through the moments of
degree at most $9n$, so $\rank\mathsf F_n\le C(n+1)^3$. The endpoint
bounds $\abs{\mathsf p_N(z)-1}\le(4r(1+r))^N/(2(1-2r))$ for
$\abs z\le r$ and $\abs{\mathsf p_N(z)}\le(4r(1+r))^N$ for
$\abs{z-1}\le r$, $r<\frac12$, together with the Taylor tail in
six variables,
\[
\sum_{k\ge n}\binom{k+5}5\Bigl(\frac r{r_0}\Bigr)^k
\le C(n+1)^5\Bigl(\frac r{r_0}\Bigr)^n,\qquad r\le\frac{r_0}2,
\]
applied at $r=3/s$, show
that on the radius-$3/s$ complex neighbourhoods of the four blocks
$B_2(\cdot)\times B_2(\cdot)$ of $(B_2(0)\cup B_2(s\eone))^2$, the
target kernel $\mathfrak K_\bullet$ of $\JointD_{\mu,s}$, equal to
$-\mathfrak X_\mu$ (for $\bullet=f$) and $-\mathfrak P_\mu$ (for
$\bullet=g$) on the two cross blocks $B_2(0)\times B_2(s\eone)$ and
$B_2(s\eone)\times B_2(0)$, and to zero on the two diagonal blocks,
satisfies
$\sup\abs{\mathfrak K_\bullet-q_{n,\bullet}}\le C(n+1)^6M_\bullet(c_0/s)^n$
with a universal $c_0\ge3/r_0$. The four blocks separate as
follows: on its own cross block, a summand of $q_{n,\bullet}$ differs
from the target by the Taylor tail plus the selector defects
$\mathsf p_{2n}-1$; on each of the other three blocks, where the
nonzero target (if any) is matched by the other summand, the evaluated
Taylor polynomial is finite with at most geometric coefficient growth,
$\abs{\mathcal T_{n,\bullet}}\le C(n+1)^5M_\bullet(2/r_0)^n$ there,
while at least one of its selector factors sits at its order-$2n$ zero
and is bounded by $(4r(1+r))^{2n}\le(C/s)^{2n}$, which absorbs the
growth after enlarging $c_0$; no convergence of the Taylor series
outside its polydisc is used.
Then Cauchy estimates bound the derivatives of order at most two on the
real blocks by $CM_\bullet(c_{\mathrm{pol}}/s)^n$ with
$c_{\mathrm{pol}}=2^6c_0$. By \eqref{eq:smooth-kernel},
$\norm{\JointD_{\mu,s}-\mathsf F_n}\le Cs^{-2}(1+\xi)^2e^{-\xi/2}(c_{\mathrm{pol}}/s)^n$,
which gives the convergence with
$S_{\mathrm{pol}}=\max\{5,2c_{\mathrm{pol}}\}$ (the bound
\eqref{eq:smooth-kernel} requires $s\ge5$; and since $c_0\ge3/r_0$,
$s\ge2c_{\mathrm{pol}}$ keeps the complex neighbourhoods of the
blocks, of radius $3/s$ in the scaled variables, inside the polydisc
of holomorphy), and
multiplied by the rank of $\mathsf G_n$ gives
\eqref{eq:increment-unweighted}. For \eqref{eq:increment-weighted},
expand $q_{n+1,\bullet}-q_{n,\bullet}$ in the midpoint monomials
$(({\bx-s\eone/2})/s)^\alpha(({\by-s\eone/2})/s)^\beta$; the coefficient sums
are at most $CM_\bullet K_{\mathrm{mid}}^n$ for a universal
$K_{\mathrm{mid}}$, because the selectors and shifts multiply
coefficient sums by at most geometric factors. Lemma~\ref{lem:moments}
at $b=\frac12$, with $U=U_s\subset B_T(s\eone/2)$, $T=s/2+1\le s$,
bounds each normalized moment
$s^{-\abs\alpha}\mathsf m^\bullet_\alpha$ by
$M_{1/2}\,T(1+\mu T)^{-1}e^{\Lambda_0\mu T}B_{1/2}^{\abs\alpha}$ for
$\bullet=f$, and with $T^2(1+\mu T)^{-2}$ in place of
$T(1+\mu T)^{-1}$ for $\bullet=g$, in the sense that the moment is
$\ip{\psi}{\JointRes^{1/2}\,\cdot\,}$ for a vector $\psi$ of at most
that norm, so the increment is
$\JointRes^{1/2}\mathsf B_n\JointRes^{1/2}$ with $\mathsf B_n$ a finite combination of
rank-one operators of total trace norm at most
$C(1+K_{\mathrm{mid}})B_{1/2}^9M_{1/2}^2e^{2\Lambda_0\mu T}(K_{\mathrm{mid}}B_{1/2}^9)^n$
times $[M_fT^2(1+\mu T)^{-2}+M_gT^4(1+\mu T)^{-4}]\le Ce^{-\xi/2}$,
using $T\le s$ and $1+\mu T\ge(1+\xi)/2$. This is
\eqref{eq:increment-weighted} with $K_{\mathrm{pol}}=K_{\mathrm{mid}}B_{1/2}^9$.
\end{proof}

\begin{proof}[Proof of Theorem~\ref{thm:reverse}]
Let $\theta=1/(4\Lambda_0)=2^{-17}$, so that $a=\theta/2$ and
$\gamma=2(1-\theta)$. For each increment, the analytic family
$z\mapsto\JointRes^{(1-z)/2}\mathsf B_n\JointRes^{(1-z)/2}$ on
$0\le\Rea z\le1$ is continuous in trace norm and holomorphic inside
(it is a finite combination of rank-one operators, and
$w\mapsto\JointRes^w$ is strongly continuous on $\Rea w\ge0$), with
boundary trace norms $\norm{\mathsf G_n}_1$ at $\Rea z=0$, where the
imaginary powers are unitary, and $\norm{\mathsf B_n}_1$ at $\Rea z=1$. Hadamard's
three-lines theorem for the trace norm \cite{Simon05} at $z=\theta$
gives
$Y_n=\JointRes^{(1-\theta)/2}\mathsf B_n\JointRes^{(1-\theta)/2}$ with
$\norm{Y_n}_1\le\norm{\mathsf G_n}_1^{1-\theta}\norm{\mathsf B_n}_1^\theta$ and
$\mathsf G_n=\JointRes^aY_n\JointRes^a$. Inserting
\eqref{eq:increment-unweighted} and \eqref{eq:increment-weighted}, with
$2\Lambda_0\theta=\frac12$, $s^{-2(1-\theta)}=s^{-\gamma}$ and
$\mu T/2\le3\xi/8$ for $s\ge4$,
\[
\norm{Y_n}_1\le C_1^{1-\theta}C_2^\theta s^{-\gamma}(1+\xi)^2e^{-\xi/8}
(n+1)^3\bigl[(c_{\mathrm{pol}}/s)^{1-\theta}K_{\mathrm{pol}}^\theta\bigr]^n .
\]
Choose $S_{\mathrm{rev}}\ge\max\{S_{\mathrm{pol}},5\}$ so large that
the bracket is at most $\frac12$ for $s\ge S_{\mathrm{rev}}$. Since
$\sum_{n\ge0}(n+1)^32^{-n}=52$ and $\sup_{\xi\ge0}(1+\xi)^2e^{-\xi/8}\le256$,
the series $Y_{\mu,s}=\sum_nY_n$ converges in trace norm with
$\norm{Y_{\mu,s}}_1\le\mathcal J_0s^{-\gamma}$,
$\mathcal J_0=52\cdot256\,C_1^{1-\theta}C_2^\theta$. Each $Y_n$ is
self-adjoint because $\mathsf B_n$ is, and
$\JointRes^aY_{\mu,s}\JointRes^a=\sum_n\mathsf G_n=\lim_N\mathsf F_{N+1}=\JointD_{\mu,s}$.
\end{proof}

% Reference list ragged right; no line break right after "doi:", "arXiv:"
% or "https:", and breaks after hyphens in URLs only where unavoidable.
\raggedright
\def\UrlNoBreaks{\do\(\do\[\do\{\do\<\do\:}\mathchardef\UrlBigBreakPenalty=900\relax
\bibliography{vacuum-bell}

\end{document}

